\input{preamble}

% OK, commencer ici.
%
\begin{document}

\title{Hyperrecouvrements}


\maketitle

\phantomsection
\label{section-phantom}

\tableofcontents

\section{Introduction}
\label{section-introduction}

\noindent
Soit $\mathcal{C}$ un site ; voir Sites, Définition \ref{sites-definition-site}.
Soit $X$ un objet de $\mathcal{C}$.
Étant donné un faisceau abélien $\mathcal{F}$
sur $\mathcal{C}$, nous voudrions calculer
ses groupes de cohomologie
$$
H^i(X, \mathcal{F}).
$$
D'après nos définitions générales (Cohomologie sur les sites, section
\ref{sites-cohomology-section-cohomology-sheaves})
ce groupe de cohomologie se calcule en
choisissant une résolution injective
$
0 \to \mathcal{F} \to \mathcal{I}^0 \to \mathcal{I}^1 \to \ldots
$
et en posant
$$
H^i(X, \mathcal{F})
=
H^i(
\Gamma(X, \mathcal{I}^0) \to
\Gamma(X, \mathcal{I}^1) \to
\Gamma(X, \mathcal{I}^2)\to \ldots)
$$
Le but de ce chapitre est de montrer que l'on peut aussi calculer ces
groupes de cohomologie sans choisir de résolution injective
(lorsque $\mathcal{C}$ admet des produits fibrés). Pour cela,
nous utiliserons des hyperrecouvrements.

\medskip\noindent
Un hyperrecouvrement dans un site est une généralisation d'un recouvrement ; voir
\cite[Expos\'e V, Sec. 7]{SGA4}. Étant donné un hyperrecouvrement $K$ d'un objet
$X$, il existe une suite spectrale de {\v C}ech à la cohomologie
qui exprime la cohomologie d'un faisceau abélien $\mathcal{F}$
sur $X$ en fonction de la cohomologie du faisceau sur les
composantes $K_n$ de $K$. Il se trouve qu'il existe toujours
assez d'hyperrecouvrements ; en prenant donc la limite inductive sur tous les hyperrecouvrements,
la suite spectrale dégénère et la cohomologie de $\mathcal{F}$
sur $X$ se calcule par la limite inductive des groupes de cohomologie de {\v C}ech.

\medskip\noindent
Un dispositif plus général que l'on peut considérer est une augmentation simpliciale pour laquelle
on a descente cohomologique ; voir \cite[Expos\'e Vbis]{SGA4}. Le texte de
Brian Conrad constitue un bel exposé sur la descente cohomologique ; voir
\url{https://math.stanford.edu/~conrad/papers/hypercover.pdf}.
Nous reviendrons sur ces questions dans le chapitre consacré aux espaces simpliciaux,
où nous montrerons, par exemple, que les hyperrecouvrements propres d'espaces
topologiques ``localement compacts'' sont de descente cohomologique
(Espaces simpliciaux, section
\ref{spaces-simplicial-section-proper-hypercovering}).
Notre méthode consistera à ramener cet énoncé à la suite spectrale de {\v C}ech
à la cohomologie construite dans ce chapitre.






























\section{Objets semi-représentables}
\label{section-semi-representable}

\noindent
Pour commencer, nous posons la définition suivante.
Les lettres « SR » signifient semi-représentable.

\begin{definition}
\label{definition-SR}
Soit $\mathcal{C}$ une catégorie. On note $\text{SR}(\mathcal{C})$
la catégorie des {\it objets semi-représentables}, définie comme suit :
\begin{enumerate}
\item les objets sont les familles d'objets $\{U_i\}_{i \in I}$, et
\item les morphismes $\{U_i\}_{i \in I} \to \{V_j\}_{j \in J}$ sont donnés par
une application $\alpha : I \to J$ et, pour chaque $i \in I$,
un morphisme $f_i : U_i \to V_{\alpha(i)}$ de $\mathcal{C}$.
\end{enumerate}
Soit $X \in \Ob(\mathcal{C})$ un objet de $\mathcal{C}$.
La catégorie des {\it objets semi-représentables au-dessus de $X$}
est la catégorie
$\text{SR}(\mathcal{C}, X) = \text{SR}(\mathcal{C}/X)$.
\end{definition}

\noindent
Cette définition est essentiellement équivalente à celle de
\cite[Expos\'e V, sous-section 7.3.0]{SGA4}. Notons qu'il s'agit
d'une « grande » catégorie. Nous bornerons plus loin la taille des ensembles
d'indices $I$ dont nous avons besoin pour les hyperrecouvrements de $X$. Nous pourrons alors redéfinir
$\text{SR}(\mathcal{C}, X)$ de manière à obtenir une catégorie. Décrivons explicitement
les objets et les morphismes de $\text{SR}(\mathcal{C}, X)$ :
\begin{enumerate}
\item les objets sont les familles de morphismes
$\{U_i \to X\}_{i \in I}$, et
\item les morphismes $\{U_i \to X\}_{i \in I} \to
\{V_j \to X\}_{j \in J}$ sont donnés par
une application $\alpha : I \to J$ et, pour chaque $i \in I$,
un morphisme $f_i : U_i \to V_{\alpha(i)}$ au-dessus de $X$.
\end{enumerate}
Il existe un foncteur d'oubli
$\text{SR}(\mathcal{C}, X) \to \text{SR}(\mathcal{C})$.

\begin{definition}
\label{definition-SR-F}
Soit $\mathcal{C}$ une catégorie.
On note $F$ le foncteur {\it qui associe un préfaisceau à un
objet semi-représentable}. En formule,
\begin{eqnarray*}
F : \text{SR}(\mathcal{C}) & \longrightarrow & \textit{PSh}(\mathcal{C}) \\
\{U_i\}_{i \in I} & \longmapsto & \amalg_{i\in I} h_{U_i}
\end{eqnarray*}
où $h_U$ désigne le préfaisceau représentable associé à
l'objet $U$.
\end{definition}

\noindent
À un morphisme $U \to X$ est associé un morphisme $h_U \to h_X$ de préfaisceaux
représentables. On considère donc souvent $F$ sur $\text{SR}(\mathcal{C}, X)$
comme un foncteur à valeurs dans la catégorie des préfaisceaux d'ensembles au-dessus de $h_X$,
à savoir $\textit{PSh}(\mathcal{C})/h_X$. Voici le diagramme correspondant :
$$
\xymatrix{
\text{SR}(\mathcal{C}, X) \ar[r]_F \ar[d] &
\textit{PSh}(\mathcal{C})/h_X \ar[d] \\
\text{SR}(\mathcal{C}) \ar[r]^F &
\textit{PSh}(\mathcal{C})
}
$$
Étudions maintenant l'existence de limites dans la catégorie des objets semi-
représentables.

\begin{lemma}
\label{lemma-coprod-prod-SR}
Soit $\mathcal{C}$ une catégorie.
\begin{enumerate}
\item la catégorie $\text{SR}(\mathcal{C})$ admet des coproduits
et $F$ commute à ceux-ci,
\item le foncteur $F : \text{SR}(\mathcal{C}) \to \textit{PSh}(\mathcal{C})$
commute aux limites,
\item si $\mathcal{C}$ admet des produits fibrés, alors $\text{SR}(\mathcal{C})$
admet des produits fibrés,
\item si $\mathcal{C}$ admet les produits de deux objets, alors
$\text{SR}(\mathcal{C})$ admet les produits de deux objets,
\item si $\mathcal{C}$ admet des égalisateurs, il en est de même de $\text{SR}(\mathcal{C})$, et
\item si $\mathcal{C}$ admet un objet final, il en est de même de $\text{SR}(\mathcal{C})$.
\end{enumerate}
Soit $X \in \Ob(\mathcal{C})$.
\begin{enumerate}
\item la catégorie $\text{SR}(\mathcal{C}, X)$ admet des coproduits
et $F$ commute à ceux-ci,
\item si $\mathcal{C}$ admet des produits fibrés, alors $\text{SR}(\mathcal{C}, X)$
admet des limites finies et
$F : \text{SR}(\mathcal{C}, X) \to \textit{PSh}(\mathcal{C})/h_X$
commute à celles-ci.
\end{enumerate}
\end{lemma}

\begin{proof}
Démonstration des assertions sur $\text{SR}(\mathcal{C})$.
Démonstration de (1). Le coproduit de $\{U_i\}_{i \in I}$ et $\{V_j\}_{j \in J}$ est
$\{U_i\}_{i \in I} \amalg \{V_j\}_{j \in J}$, autrement dit la famille
d'objets dont l'ensemble d'indices est $I \amalg J$ et qui, pour un élément
$k \in I \amalg J$, donne $U_i$ si $k = i \in I$ et donne $V_j$ si
$k = j \in J$. Il en va de même pour les coproduits
de familles d'objets. Il est clair que $F$ commute à ceux-ci.

\medskip\noindent
Démonstration de (2). Pour $U$ dans $\Ob(\mathcal{C})$, considérons l'objet $\{U\}$ de
$\text{SR}(\mathcal{C})$. Il est clair que
$\Mor_{\text{SR}(\mathcal{C})}(\{U\}, K)) = F(K)(U)$
pour $K \in \Ob(\text{SR}(\mathcal{C}))$. Comme les limites de préfaisceaux
se calculent au niveau des sections
(Sites, section \ref{sites-section-limits-colimits-PSh}),
on conclut que $F$ commute aux limites.

\medskip\noindent
Démonstration de (3). Supposons donnés un morphisme
$(\alpha, f_i) : \{U_i\}_{i \in I} \to \{V_j\}_{j \in J}$
et un morphisme
$(\beta, g_k) : \{W_k\}_{k \in K} \to \{V_j\}_{j \in J}$.
Le produit fibré de ces morphismes est donné par
$$
\{ U_i \times_{f_i, V_j, g_k} W_k\}_{(i, j, k) \in I \times J \times K
\text{ tels que } j = \alpha(i) = \beta(k)}
$$
Les produits fibrés existent si $\mathcal{C}$ admet des produits fibrés.

\medskip\noindent
Démonstration de (4). Le produit de $\{U_i\}_{i \in I}$ et $\{V_j\}_{j \in J}$ est
$\{U_i \times V_j\}_{i \in I, j \in J}$. Les produits existent si
$\mathcal{C}$ admet des produits.

\medskip\noindent
Démonstration de (5). L'égalisateur de deux applications
$(\alpha, f_i), (\alpha', f'_i) : \{U_i\}_{i \in I} \to \{V_j\}_{j \in J}$
est
$$
\{
\text{Eq}(f_i, f'_i : U_i \to V_{\alpha(i)})
\}_{i \in I,\ \alpha(i) = \alpha'(i)}
$$
Les égalisateurs existent si $\mathcal{C}$ admet des égalisateurs.

\medskip\noindent
Démonstration de (6). Si $X$ est un objet final de $\mathcal{C}$, alors
$\{X\}$ est un objet final de $\text{SR}(\mathcal{C})$.

\medskip\noindent
Démonstration des assertions sur $\text{SR}(\mathcal{C}, X)$.
Elles résultent des résultats précédents appliqués à la catégorie
$\mathcal{C}/X$, compte tenu de
$\text{SR}(\mathcal{C}/X) = \text{SR}(\mathcal{C}, X)$ et de
$\textit{PSh}(\mathcal{C}/X) = \textit{PSh}(\mathcal{C})/h_X$
(Sites, Lemme \ref{sites-lemma-essential-image-j-shriek} appliqué
à $\mathcal{C}$ munie de la topologie chaotique). Toutefois,
donnons aussi l'argument direct suivant.
Il est clair que le coproduit de
$\{U_i \to X\}_{i \in I}$ et $\{V_j \to X\}_{j \in J}$
est $\{U_i \to X\}_{i \in I} \amalg \{V_j \to X\}_{j \in J}$,
et de même pour les coproduits de
familles de familles de morphismes de but $X$.
L'objet $\{X \to X\}$ est un objet final
de $\text{SR}(\mathcal{C}, X)$.
Supposons donnés un morphisme
$(\alpha, f_i) : \{U_i \to X\}_{i \in I} \to \{V_j \to X\}_{j \in J}$
et un morphisme
$(\beta, g_k) : \{W_k \to X\}_{k \in K} \to \{V_j \to X\}_{j \in J}$.
Le produit fibré de ces morphismes est donné par
$$
\{ U_i \times_{f_i, V_j, g_k} W_k \to X \}_{(i, j, k) \in I \times J \times K
\text{ tels que } j = \alpha(i) = \beta(k)}
$$
Les produits fibrés existent d'après l'hypothèse selon laquelle
$\mathcal{C}$ admet des produits fibrés.
Ainsi $\text{SR}(\mathcal{C}, X)$ admet des limites finies ;
voir Catégories, Lemme \ref{categories-lemma-finite-limits-exist}.
Nous omettons la vérification des assertions sur le foncteur $F$ dans ce cas.
\end{proof}





\section{Hyperrecouvrements}
\label{section-hypercoverings}

\noindent
Si nous supposons que notre catégorie est un site, nous pouvons poser la
définition suivante.

\begin{definition}
\label{definition-covering-SR}
Soit $\mathcal{C}$ un site. Soit
$f = (\alpha, f_i) : \{U_i\}_{i \in I} \to \{V_j\}_{j \in J}$
un morphisme de la catégorie $\text{SR}(\mathcal{C})$.
On dit que $f$ est un {\it recouvrement} si, pour tout $j \in J$, la
famille de morphismes $\{U_i \to V_j\}_{i \in I, \alpha(i) = j}$
est couvrante pour le site $\mathcal{C}$.
Soit $X$ un objet de $\mathcal{C}$.
Un morphisme $K \to L$ de $\text{SR}(\mathcal{C}, X)$ est
un {\it recouvrement} si son image dans $\text{SR}(\mathcal{C})$ est
un recouvrement.
\end{definition}

\begin{lemma}
\label{lemma-covering-permanence}
Soit $\mathcal{C}$ un site.
\begin{enumerate}
\item Un composé de recouvrements dans $\text{SR}(\mathcal{C})$
est un recouvrement.
\item Si $K \to L$ est un recouvrement dans $\text{SR}(\mathcal{C})$
et $L' \to L$ est un morphisme, alors $L' \times_L K$ existe
et $L' \times_L K \to L'$ est un recouvrement.
\item Si $\mathcal{C}$ admet les produits de deux objets, et si
$A \to B$ et $K \to L$ sont des recouvrements dans $\text{SR}(\mathcal{C})$,
alors $A \times K \to B \times L$ est un recouvrement.
\end{enumerate}
Soit $X \in \Ob(\mathcal{C})$. Alors (1) et (2) valent pour
$\text{SR}(\mathcal{C}, X)$, et (3) vaut si $\mathcal{C}$
admet des produits fibrés.
\end{lemma}

\begin{proof}
L'assertion (1) résulte immédiatement des axiomes d'un site.
L'assertion (2) résulte de la construction des produits fibrés
dans $\text{SR}(\mathcal{C})$ donnée dans la démonstration du
Lemme \ref{lemma-coprod-prod-SR}
et de l'exigence que les morphismes d'un recouvrement
de $\mathcal{C}$ soient représentables.
L'assertion (3) s'obtient en considérant $A \times K \to B \times L$
comme le composé $A \times K \to B \times K \to B \times L$,
donc comme un composé de changements de base de recouvrements.
La dernière assertion résulte de l'égalité $\text{SR}(\mathcal{C}, X) =
\text{SR}(\mathcal{C}/X)$.
\end{proof}

\noindent
D'après le Lemme \ref{lemma-coprod-prod-SR} et le
Lemme \ref{simplicial-lemma-existence-cosk} de Simplicial,
le cosquelette d'un objet simplicial tronqué de
$\text{SR}(\mathcal{C}, X)$ existe si $\mathcal{C}$ admet des produits fibrés.
La définition suivante a donc un sens.

\begin{definition}
\label{definition-hypercovering}
Soit $\mathcal{C}$ un site. Supposons que $\mathcal{C}$ admette des produits fibrés.
Soit $X \in \Ob(\mathcal{C})$ un objet de $\mathcal{C}$.
Un {\it hyperrecouvrement de $X$} est un objet simplicial
$K$ de $\text{SR}(\mathcal{C}, X)$ tel que
\begin{enumerate}
\item L'objet $K_0$ est un recouvrement de $X$ pour le site $\mathcal{C}$.
\item Pour tout $n \geq 0$, le morphisme canonique
$$
K_{n + 1} \longrightarrow (\text{cosk}_n \text{sk}_n K)_{n + 1}
$$
est un recouvrement au sens défini ci-dessus.
\end{enumerate}
\end{definition}

\noindent
La condition (1) a un sens puisque tout objet de
$\text{SR}(\mathcal{C}, X)$ est, après tout, une famille
de morphismes de but $X$. On pourrait aussi la
formuler en disant que le morphisme de $K_0$ vers
l'objet final de $\text{SR}(\mathcal{C}, X)$
est un recouvrement.

\begin{example}[Hyperrecouvrements de {\v C}ech]
\label{example-cech}
Soit $\mathcal{C}$ un site admettant des produits fibrés.
Soit $\{U_i \to X\}_{i \in I}$ un recouvrement de $\mathcal{C}$.
Posons $K_0 = \{U_i \to X\}_{i \in I}$.
Alors $K_0$ est un objet simplicial $0$-tronqué de
$\text{SR}(\mathcal{C}, X)$. Nous pouvons donc former
$$
K = \text{cosk}_0 K_0.
$$
Il est clair que $K$ vérifie la condition (1) de la Définition \ref{definition-hypercovering}.
Comme tous les morphismes $K_{n + 1} \to (\text{cosk}_n \text{sk}_n K)_{n + 1}$
sont des isomorphismes d'après le
Lemme \ref{simplicial-lemma-cosk-up} de Simplicial,
il vérifie aussi la condition (2). Remarquons que
les termes $K_n$ sont les objets usuels
$$
K_n = \{
U_{i_0} \times_X U_{i_1} \times_X \ldots \times_X U_{i_n} \to X
\}_{(i_0, i_1, \ldots, i_n) \in I^{n + 1}}
$$
Un hyperrecouvrement de $X$ de cette forme est appelé un
{\it hyperrecouvrement de {\v C}ech} de $X$.
\end{example}

\begin{example}[Hyperrecouvrement par un objet simplicial du site]
\label{example-hypercovering-in-C}
Soit $\mathcal{C}$ un site admettant des produits fibrés. Soit
$X \in \Ob(\mathcal{C})$. Soit $U$ un objet simplicial de $\mathcal{C}$.
Comme d'habitude, on note $U_n = U([n])$. Enfin, supposons donnée une augmentation
$$
a : U \to X
$$
Dans cette situation, on peut considérer l'objet simplicial $K$
de $\text{SR}(\mathcal{C}, X)$ dont les termes sont $K_n = \{U_n \to X\}$.
Alors $K$ est un hyperrecouvrement de $X$ au sens de la
Définition \ref{definition-hypercovering}
si et seulement si les trois
conditions suivantes\footnote{Comme $\mathcal{C}$ admet des produits fibrés, la
catégorie $\mathcal{C}/X$ admet toutes les limites finies.
Les cosquelettes requis existent donc d'après le
Lemme \ref{simplicial-lemma-existence-cosk} de Simplicial.} sont vérifiées :
\begin{enumerate}
\item $\{U_0 \to X\}$ est un recouvrement de $\mathcal{C}$,
\item $\{U_1 \to U_0 \times_X U_0\}$ est un recouvrement de $\mathcal{C}$,
\item $\{U_{n + 1} \to (\text{cosk}_n\text{sk}_n U)_{n + 1}\}$
est un recouvrement de $\mathcal{C}$ pour $n \geq 1$.
\end{enumerate}
Nous omettons la vérification immédiate.
\end{example}

\begin{example}[Hyperrecouvrement de {\v C}ech associé à un recouvrement]
\label{example-cech-cover}
Soit $\mathcal{C}$ un site admettant des produits fibrés. Soit $U \to X$ un
morphisme de $\mathcal{C}$ tel que $\{U \to X\}$ soit un recouvrement de
$\mathcal{C}$\footnote{Un morphisme de $\mathcal{C}$ ayant cette propriété
est parfois appelé un « recouvrement ».}. Considérons l'objet simplicial $K$ de
$\text{SR}(\mathcal{C}, X)$ dont les termes sont
$$
K_n = \{U \times_X U \times_X \ldots \times_X U \to X\}
\quad (n + 1 \text{ facteurs})
$$
Alors $K$ est un hyperrecouvrement de $X$. Cet exemple est un cas particulier à la fois de
l'Exemple \ref{example-cech} et de
l'Exemple \ref{example-hypercovering-in-C}.
\end{example}

\begin{lemma}
\label{lemma-hypercoverings-set}
Soit $\mathcal{C}$ un site admettant des produits fibrés.
Soit $X \in \Ob(\mathcal{C})$ un objet de $\mathcal{C}$.
La collection de tous les hyperrecouvrements de $X$ forme un ensemble.
\end{lemma}

\begin{proof}
Puisque $\mathcal{C}$ est un site, l'ensemble de tous les recouvrements de
$X$ forme un ensemble. On voit donc que la collection
des $K_0$ possibles forme un ensemble. Supposons avoir montré que
la collection de tous les $K_0, \ldots, K_n$ possibles forme
un ensemble. Il suffit alors de montrer qu'étant donnés
$K_0, \ldots, K_n$, la collection de tous les
$K_{n + 1}$ possibles forme un ensemble. Cela est clair, puisque
il faut choisir $K_{n + 1}$ parmi tous les recouvrements possibles
de $(\text{cosk}_n \text{sk}_n K)_{n + 1}$.
\end{proof}

\begin{remark}
\label{remark-hypercoverings-really-set}
Le lemme ne dit pas seulement qu'il existe un système cofinal
de choix d'hyperrecouvrements qui soit un ensemble,
mais que les hyperrecouvrements eux-mêmes forment effectivement un ensemble.
\end{remark}

\noindent
La catégorie des préfaisceaux sur $\mathcal{C}$ admet des
(co)limites finies. Les foncteurs $\text{cosk}_n$
existent donc pour les préfaisceaux d'ensembles.

\begin{lemma}
\label{lemma-hypercovering-F}
Soit $\mathcal{C}$ un site admettant des produits fibrés.
Soit $X \in \Ob(\mathcal{C})$ un objet de $\mathcal{C}$.
Soit $K$ un hyperrecouvrement de $X$.
Considérons l'objet simplicial $F(K)$ de $\textit{PSh}(\mathcal{C})$,
muni de son augmentation vers le préfaisceau simplicial constant $h_X$.
\begin{enumerate}
\item Le morphisme de préfaisceaux $F(K)_0 \to h_X$ devient
surjectif après passage aux faisceaux associés.
\item Le morphisme
$$
(d^1_0, d^1_1) :
F(K)_1
\longrightarrow
F(K)_0 \times_{h_X} F(K)_0
$$
devient surjectif après passage aux faisceaux associés.
\item Pour tout $n \geq 1$, le morphisme
$$
F(K)_{n + 1} \longrightarrow (\text{cosk}_n \text{sk}_n F(K))_{n + 1}
$$
devient surjectif après passage aux faisceaux associés.
\end{enumerate}
\end{lemma}

\begin{proof}
Nous utiliserons le fait que si
$\{U_i \to U\}_{i \in I}$ est un recouvrement du site
$\mathcal{C}$, alors le morphisme
$$
\amalg_{i \in I} h_{U_i} \to h_U
$$
devient surjectif après passage aux faisceaux associés ; voir
Sites, Lemme \ref{sites-lemma-covering-surjective-after-sheafification}.
La première assertion en résulte immédiatement.

\medskip\noindent
Pour la deuxième assertion, remarquons que, d'après
Simplicial, Exemple \ref{simplicial-example-cosk0},
l'objet simplicial $\text{cosk}_0 \text{sk}_0 K$
a pour termes $K_0 \times \ldots \times K_0$. Ainsi,
d'après la définition d'un hyperrecouvrement,
on voit que $(d^1_0, d^1_1) : K_1 \to K_0 \times K_0$ est un
recouvrement. L'assertion (2) résulte donc de l'observation précédente
et du fait que $F$ transforme les produits en produits
fibrés au-dessus de $h_X$.

\medskip\noindent
Pour la troisième assertion, nous affirmons que
$\text{cosk}_n \text{sk}_n F(K) =
F(\text{cosk}_n \text{sk}_n K)$ pour $n \geq 1$.
Pour le montrer, notons temporairement $F'$ le foncteur
$\text{SR}(\mathcal{C}, X) \to \textit{PSh}(\mathcal{C})/h_X$.
D'après le Lemme \ref{lemma-coprod-prod-SR}, le foncteur
$F'$ commute aux limites finies.
La description du foncteur $\text{cosk}_n$ donnée dans
Simplicial, section \ref{simplicial-section-skeleton},
montre que $\text{cosk}_n \text{sk}_n F'(K) =
F'(\text{cosk}_n \text{sk}_n K)$.
Rappelons que la catégorie intervenant dans la description de
$(\text{cosk}_n U)_m$ du
Lemme \ref{simplicial-lemma-existence-cosk} de Simplicial
est la catégorie $(\Delta/[m])^{opp}_{\leq n}$. C'est un
exercice plaisant que de montrer que $(\Delta/[m])_{\leq n}$ est
une catégorie connexe (voir
Catégories, Définition \ref{categories-definition-category-connected})
dès que $n \geq 1$. Par conséquent, le
Lemme \ref{categories-lemma-connected-limit-over-X} de Catégories
montre que $\text{cosk}_n \text{sk}_n F'(K) =
\text{cosk}_n \text{sk}_n F(K)$. D'où l'assertion.
L'assertion (3) en résulte, car on voit maintenant que
le morphisme de (3) s'obtient en appliquant le
foncteur $F$ à un recouvrement comme dans la Définition \ref{definition-covering-SR},
et le résultat découle du premier fait mentionné
dans cette démonstration.
\end{proof}



\section{Acyclicité}
\label{section-acyclicity}

\noindent
Soit $\mathcal{C}$ un site.
Pour un préfaisceau d'ensembles $\mathcal{F}$, on note $\mathbf{Z}_\mathcal{F}$
le préfaisceau de groupes abéliens défini par la règle
$$
\mathbf{Z}_\mathcal{F}(U) = \text{groupe abélien libre sur }\mathcal{F}(U).
$$
Nous l'appellerons parfois le {\it préfaisceau abélien libre sur $\mathcal{F}$}.
Bien entendu, la construction $\mathcal{F} \mapsto \mathbf{Z}_\mathcal{F}$
est un foncteur, adjoint à gauche du foncteur d'oubli
$\textit{PAb}(\mathcal{C}) \to \textit{PSh}(\mathcal{C})$.
Bien entendu, le faisceau associé $\mathbf{Z}_\mathcal{F}^\#$ est
un faisceau de groupes abéliens, et le foncteur
$\mathcal{F} \mapsto \mathbf{Z}_\mathcal{F}^\#$ est lui aussi
un adjoint à gauche. Nous appelons parfois $\mathbf{Z}_\mathcal{F}^\#$
le {\it faisceau abélien libre sur $\mathcal{F}$}.

\medskip\noindent
Pour un objet $X$ du site $\mathcal{C}$, on note
$\mathbf{Z}_X$ le préfaisceau abélien libre sur $h_X$, et
on note $\mathbf{Z}_X^\#$ son faisceau associé.

\begin{definition}
\label{definition-homology}
Soit $\mathcal{C}$ un site.
Soit $K$ un objet simplicial de $\textit{PSh}(\mathcal{C})$.
Par ce qui précède, on obtient un objet simplicial $\mathbf{Z}_K^\#$ de
$\textit{Ab}(\mathcal{C})$. On peut prendre son
complexe associé de préfaisceaux abéliens $s(\mathbf{Z}_K^\#)$ ; voir
Simplicial, section \ref{simplicial-section-complexes}.
L'{\it homologie de $K$} est l'homologie du
complexe de faisceaux abéliens $s(\mathbf{Z}_K^\#)$.
\end{definition}

\noindent
Autrement dit, la {\it $i$-ième homologie $H_i(K)$ de $K$}
est le faisceau de groupes abéliens $H_i(K) = H_i(s(\mathbf{Z}_K^\#))$.
Dans cette section, nous étudions l'homologie lorsque $K$
est un hyperrecouvrement d'un objet $X$ de $\mathcal{C}$.

\begin{lemma}
\label{lemma-compare-cosk0}
Soit $\mathcal{C}$ un site.
Soit $\mathcal{F} \to \mathcal{G}$ un morphisme
de préfaisceaux d'ensembles. Notons $K$ l'objet
simplicial de $\textit{PSh}(\mathcal{C})$ dont le $n$-ième
terme est le $(n + 1)$-ième produit fibré de $\mathcal{F}$
au-dessus de $\mathcal{G}$ ; voir
Simplicial, Exemple \ref{simplicial-example-fibre-products-simplicial-object}.
Alors, si $\mathcal{F} \to \mathcal{G}$ devient surjectif après
passage aux faisceaux associés, on a
$$
H_i(K) =
\left\{
\begin{matrix}
0 & \text{si} & i > 0\\
\mathbf{Z}_\mathcal{G}^\# & \text{si} & i = 0
\end{matrix}
\right.
$$
L'isomorphisme en degré $0$ est donné par le
morphisme $H_0(K) \to \mathbf{Z}_\mathcal{G}^\#$
provenant de l'application $(\mathbf{Z}_K^\#)_0 =
\mathbf{Z}_\mathcal{F}^\# \to \mathbf{Z}_\mathcal{G}^\#$.
\end{lemma}

\begin{proof}
Soit $\mathcal{G}' \subset \mathcal{G}$ l'image du
morphisme $\mathcal{F} \to \mathcal{G}$.
Soit $U \in \Ob(\mathcal{C})$. Posons
$A = \mathcal{F}(U)$ et $B = \mathcal{G}'(U)$.
Alors l'ensemble simplicial $K(U)$ est égal à l'ensemble
simplicial dont les $n$-simplexes sont donnés par
$$
A \times_B A \times_B \ldots \times_B A\ (n + 1 \text{ facteurs)}.
$$
D'après Simplicial, Lemme \ref{simplicial-lemma-cosk-minus-one-equivalence},
le morphisme $K(U) \to B$ est une fibration de Kan triviale.
C'est donc une équivalence d'homotopie
(Simplicial, Lemme \ref{simplicial-lemma-trivial-kan-homotopy}).
En appliquant le foncteur « groupe abélien libre sur » à ce morphisme,
on en déduit que
$$
\mathbf{Z}_K(U) \longrightarrow \mathbf{Z}_B
$$
est une équivalence d'homotopie. Remarquons que $s(\mathbf{Z}_B)$ est
le complexe
$$
\ldots \to
\bigoplus\nolimits_{b \in B}\mathbf{Z} \xrightarrow{0}
\bigoplus\nolimits_{b \in B}\mathbf{Z} \xrightarrow{1}
\bigoplus\nolimits_{b \in B}\mathbf{Z} \xrightarrow{0}
\bigoplus\nolimits_{b \in B}\mathbf{Z} \to 0
$$
voir Simplicial, Lemme \ref{simplicial-lemma-homology-eilenberg-maclane}.
On voit donc que
$H_i(s(\mathbf{Z}_K(U))) = 0$ pour $i > 0$, et
$H_0(s(\mathbf{Z}_K(U))) = \bigoplus_{b \in B}\mathbf{Z}
= \bigoplus_{s \in \mathcal{G}'(U)} \mathbf{Z}$.
Ces identifications sont compatibles aux morphismes de
restriction.

\medskip\noindent
On conclut que $H_i(s(\mathbf{Z}_K)) = 0$ pour $i > 0$ et que
$H_0(s(\mathbf{Z}_K)) = \mathbf{Z}_{\mathcal{G}'}$, où l'on
calcule ici les groupes d'homologie dans $\textit{PAb}(\mathcal{C})$. Comme
le foncteur « faisceau associé » est exact, on en déduit le résultat
du lemme. En effet, l'exactitude implique
que $H_0(s(\mathbf{Z}_K))^\# = H_0(s(\mathbf{Z}_K^\#))$,
et de même pour les autres indices.
\end{proof}

\begin{lemma}
\label{lemma-acyclicity}
Soit $\mathcal{C}$ un site.
Soit $f : L \to K$ un morphisme d'objets
simpliciaux de $\textit{PSh}(\mathcal{C})$.
Soit $n \geq 0$ un entier.
Supposons que
\begin{enumerate}
\item Pour $i < n$, le morphisme $L_i \to K_i$ est un isomorphisme.
\item Le morphisme $L_n \to K_n$ devient surjectif après passage aux faisceaux associés.
\item L'application canonique $L \to \text{cosk}_n \text{sk}_n L$ est un isomorphisme.
\item L'application canonique $K \to \text{cosk}_n \text{sk}_n K$ est un isomorphisme.
\end{enumerate}
Alors $H_i(f) : H_i(L) \to H_i(K)$ est un isomorphisme.
\end{lemma}

\begin{proof}
Cette démonstration est exactement la même que celle du
Lemme \ref{lemma-compare-cosk0} ci-dessus. En effet,
on commence par noter $K_n' \subset K_n$ le sous-préfaisceau
qui est l'image de l'application $L_n \to K_n$. L'hypothèse
(2) signifie que le faisceau associé à $K_n'$ est égal au
faisceau associé à $K_n$. De plus, puisque $L_i = K_i$
pour tout $i < n$, on voit que l'on obtient un préfaisceau simplicial
$n$-tronqué $U$ en prenant
$U_0 = L_0 = K_0, \ldots, U_{n - 1} = L_{n - 1} = K_{n - 1}, U_n = K'_n$.
Notons $K' = \text{cosk}_n U$, qui est un préfaisceau simplicial.
Comme on peut construire $K'_m$ par une limite finie, et
que le foncteur « faisceau associé » est exact, on voit que
$(K'_m)^\# = K_m$. Autrement dit, $(K')^\# = K^\#$.
On conclut, en utilisant encore une fois l'exactitude du faisceau associé,
que $H_i(K) = H_i(K')$. Il suffit donc de démontrer le lemme
pour le morphisme $L \to K'$ ; autrement dit, on peut
supposer que $L_n \to K_n$ est un morphisme surjectif
de {\it préfaisceaux} !

\medskip\noindent
Dans ce cas, pour tout objet $U$ de $\mathcal{C}$,
on voit que le morphisme d'ensembles simpliciaux
$$
L(U) \longrightarrow K(U)
$$
satisfait toutes les hypothèses du
Lemme \ref{simplicial-lemma-section} de Simplicial.
C'est donc une fibration de Kan triviale. En particulier, c'est
une équivalence d'homotopie
(Simplicial, Lemme \ref{simplicial-lemma-trivial-kan-homotopy}).
Ainsi
$$
\mathbf{Z}_L(U) \longrightarrow \mathbf{Z}_K(U)
$$
est elle aussi une équivalence d'homotopie, et ce pour tout $U$.
Le résultat en découle.
\end{proof}

\begin{lemma}
\label{lemma-acyclic-hypercover-sheaves}
Soit $\mathcal{C}$ un site.
Soit $K$ un préfaisceau simplicial.
Soit $\mathcal{G}$ un préfaisceau.
Soit $K \to \mathcal{G}$ une augmentation de $K$
vers $\mathcal{G}$. Supposons que
\begin{enumerate}
\item Le morphisme de préfaisceaux $K_0 \to \mathcal{G}$ devienne
surjectif après passage aux faisceaux associés.
\item Le morphisme
$$
(d^1_0, d^1_1) :
K_1
\longrightarrow
K_0 \times_\mathcal{G} K_0
$$
devienne surjectif après passage aux faisceaux associés.
\item Pour tout $n \geq 1$, le morphisme
$$
K_{n + 1} \longrightarrow (\text{cosk}_n \text{sk}_n K)_{n + 1}
$$
devienne surjectif après passage aux faisceaux associés.
\end{enumerate}
Alors $H_i(K) = 0$ pour $i > 0$ et
$H_0(K) = \mathbf{Z}_\mathcal{G}^\#$.
\end{lemma}

\begin{proof}
Notons $K^n = \text{cosk}_n \text{sk}_n K$ pour $n \geq 1$.
Définissons $K^0$ comme l'objet simplicial dont les termes
$(K^0)_n$ sont égaux au produit fibré itéré $(n + 1)$ fois
$K_0 \times_\mathcal{G} \ldots \times_\mathcal{G} K_0$ ;
voir Simplicial,
Exemple \ref{simplicial-example-fibre-products-simplicial-object}.
On a des morphismes
$$
K \longrightarrow \ldots \to K^n \to K^{n - 1} \to \ldots \to K^1 \to K^0.
$$
Les morphismes $K \to K^i$, $K^j \to K^i$ pour $j \geq i \geq 1$ proviennent
des propriétés universelles des foncteurs $\text{cosk}_n$.
Le morphisme $K^1 \to K^0$ est le morphisme canonique
de
Simplicial, Remarque \ref{simplicial-remark-augmentation}.
Rappelons aussi que $K^0 \to \text{cosk}_1 \text{sk}_1 K^0$
est un isomorphisme ; voir
Simplicial, Lemme \ref{simplicial-lemma-cosk-minus-one}.

\medskip\noindent
D'après le Lemme \ref{lemma-compare-cosk0}, on voit que
$H_i(K^0) = 0$ pour $i > 0$ et $H_0(K^0) = \mathbf{Z}_\mathcal{G}^\#$.

\medskip\noindent
Fixons $n \geq 1$. Considérons le morphisme $K^n \to K^{n - 1}$.
C'est un isomorphisme sur les termes de degré $< n$.
Remarquons que $K^n \to \text{cosk}_n \text{sk}_n K^n$ et
$K^{n - 1} \to \text{cosk}_n \text{sk}_n K^{n - 1}$
sont des isomorphismes. Remarquons que $(K^n)_n = K_n$ et
que $(K^{n - 1})_n = (\text{cosk}_{n - 1} \text{sk}_{n - 1} K)_n$.
D'après l'hypothèse, on a donc que $(K^n)_n \to (K^{n - 1})_n$
est un morphisme de préfaisceaux qui devient surjectif après
passage aux faisceaux associés. Le Lemme \ref{lemma-acyclicity} donne alors
$H_i(K^n) = H_i(K^{n - 1})$.
Combiné avec ce qui précède, ceci démontre le lemme.
\end{proof}

\begin{lemma}
\label{lemma-hypercovering-acyclic}
Soit $\mathcal{C}$ un site admettant des produits fibrés.
Soit $X$ un objet de $\mathcal{C}$.
Soit $K$ un hyperrecouvrement de $X$.
L'homologie du préfaisceau simplicial $F(K)$ est
$0$ dans les degrés $> 0$ et égale à $\mathbf{Z}_X^\#$
en degré $0$.
\end{lemma}

\begin{proof}
Combiner les Lemmes \ref{lemma-acyclic-hypercover-sheaves}
et \ref{lemma-hypercovering-F}.
\end{proof}












\section{Cohomologie de {\v C}ech et hyperrecouvrements}
\label{section-hyper-cech}

\noindent
Soit $\mathcal{C}$ un site. Considérons un préfaisceau de
groupes abéliens $\mathcal{F}$ sur le site $\mathcal{C}$.
Il définit un foncteur
\begin{eqnarray*}
\mathcal{F} : \text{SR}(\mathcal{C})^{opp}
& \longrightarrow &
\textit{Ab} \\
\{U_i\}_{i \in I} &
\longmapsto &
\prod\nolimits_{i \in I} \mathcal{F}(U_i)
\end{eqnarray*}
Ainsi, un objet simplicial $K$ de $\text{SR}(\mathcal{C})$
donne un objet cosimplicial $\mathcal{F}(K)$ de $\textit{Ab}$.
Le complexe de cochaînes $s(\mathcal{F}(K))$ associé à $\mathcal{F}(K)$
(Simplicial, section
\ref{simplicial-section-dold-kan-cosimplicial})
est appelé le complexe de {\v C}ech de $\mathcal{F}$ relativement
à l'objet simplicial $K$. On pose
$$
\check{H}^i(K, \mathcal{F})
=
H^i(s(\mathcal{F}(K))).
$$
et on l'appelle le $i$-ième groupe de cohomologie de {\v C}ech
de $\mathcal{F}$ relativement à $K$.
Dans cette section, nous démontrons des analogues de certains résultats sur la
cohomologie de {\v C}ech des recouvrements ouverts établis dans
Cohomologie, sections \ref{cohomology-section-cech},
\ref{cohomology-section-cech-functor} et
\ref{cohomology-section-cech-cohomology-cohomology}.

\begin{lemma}
\label{lemma-h0-cech}
Soit $\mathcal{C}$ un site admettant des produits fibrés.
Soit $X$ un objet de $\mathcal{C}$.
Soit $K$ un hyperrecouvrement de $X$.
Soit $\mathcal{F}$ un faisceau de groupes abéliens sur $\mathcal{C}$.
Alors $\check{H}^0(K, \mathcal{F}) = \mathcal{F}(X)$.
\end{lemma}

\begin{proof}
On a
$$
\check{H}^0(K, \mathcal{F})
=
\Ker(\mathcal{F}(K_0) \longrightarrow \mathcal{F}(K_1))
$$
Écrivons $K_0 = \{U_i \to X\}$. C'est un recouvrement dans le site
$\mathcal{C}$. De même, $K_1 \to K_0 \times K_0$
est un recouvrement dans $\text{SR}(\mathcal{C}, X)$. On peut donc
écrire $K_1 = \amalg_{i_0, i_1 \in I} \{V_{i_0i_1j} \to X\}$
de telle sorte que le morphisme $K_1 \to K_0 \times K_0$ soit donné
par les recouvrements $\{V_{i_0i_1j} \to U_{i_0} \times_X U_{i_1}\}$
du site $\mathcal{C}$. On peut alors encore identifier
$$
\check{H}^0(K, \mathcal{F})
=
\Ker(
\prod\nolimits_i \mathcal{F}(U_i)
\longrightarrow
\prod\nolimits_{i_0i_1 j} \mathcal{F}(V_{i_0i_1j})
)
$$
où l'application est évidente. La propriété de faisceau de $\mathcal{F}$
implique que $\check{H}^0(K, \mathcal{F}) = H^0(X, \mathcal{F})$.
\end{proof}

\noindent
Cette propriété caractérise bien sûr les faisceaux abéliens parmi tous les
préfaisceaux abéliens sur $\mathcal{C}$.
L'analogue de Cohomologie, Lemme \ref{lemma-injective-trivial-cech},
est ici le suivant.

\begin{lemma}
\label{lemma-injective-trivial-cech}
Soit $\mathcal{C}$ un site admettant des produits fibrés.
Soit $X$ un objet de $\mathcal{C}$.
Soit $K$ un hyperrecouvrement de $X$.
Soit $\mathcal{I}$ un faisceau injectif de groupes abéliens sur $\mathcal{C}$.
Alors
$$
\check{H}^p(K, \mathcal{I}) =
\left\{
\begin{matrix}
\mathcal{I}(X) & \text{si} & p = 0 \\
0 & \text{si} & p > 0
\end{matrix}
\right.
$$
\end{lemma}

\begin{proof}
Observons que, pour tout objet $Z = \{U_i \to X\}$ de
$\text{SR}(\mathcal{C}, X)$ et tout faisceau abélien
$\mathcal{F}$ sur $\mathcal{C}$, on a
\begin{eqnarray*}
\mathcal{F}(Z)
& = &
\prod \mathcal{F}(U_i) \\
& = &
\prod \Mor_{\textit{PSh}(\mathcal{C})}(h_{U_i}, \mathcal{F})\\
& = &
\Mor_{\textit{PSh}(\mathcal{C})}(F(Z), \mathcal{F})\\
& = &
\Mor_{\textit{PAb}(\mathcal{C})}(\mathbf{Z}_{F(Z)}, \mathcal{F}) \\
& = &
\Mor_{\textit{Ab}(\mathcal{C})}(\mathbf{Z}_{F(Z)}^\#, \mathcal{F})
\end{eqnarray*}
On voit donc que, pour tout objet simplicial $K$ de
$\text{SR}(\mathcal{C}, X)$, on a
\begin{equation}
\label{equation-identify-cech}
s(\mathcal{F}(K))
=
\Hom_{\textit{Ab}(\mathcal{C})}(s(\mathbf{Z}_{F(K)}^\#), \mathcal{F})
\end{equation}
voir la Définition \ref{definition-homology} pour les notations.
Le complexe de faisceaux $s(\mathbf{Z}_{F(K)}^\#)$ est quasi-isomorphe
à $\mathbf{Z}_X^\#$ si $K$ est un hyperrecouvrement ; voir le
Lemme \ref{lemma-hypercovering-acyclic}. On en conclut
que si $\mathcal{I}$ est un faisceau abélien injectif et si
$K$ est un hyperrecouvrement, alors le complexe $s(\mathcal{I}(K))$
est acyclique, sauf peut-être en degré $0$.
Autrement dit, on a
$$
\check{H}^i(K, \mathcal{I}) = 0
$$
pour $i > 0$. Combiné au Lemme \ref{lemma-h0-cech}, ceci démontre le lemme.
\end{proof}

\noindent
Venons-en à l'analogue de Cohomologie sur les sites, Lemme
\ref{sites-cohomology-lemma-cech-spectral-sequence}.
Soit $\mathcal{C}$ un site.
Soit $\mathcal{F}$ un faisceau de groupes abéliens sur $\mathcal{C}$.
Rappelons que $\underline{H}^i(\mathcal{F})$ désigne le préfaisceau
de groupes abéliens sur $\mathcal{C}$ défini par la
règle $\underline{H}^i(\mathcal{F}) : U \longmapsto H^i(U, \mathcal{F})$.
On l'étend à $\text{SR}(\mathcal{C})$ comme dans l'introduction
de cette section.

\begin{lemma}
\label{lemma-cech-spectral-sequence}
Soit $\mathcal{C}$ un site admettant des produits fibrés.
Soit $X$ un objet de $\mathcal{C}$.
Soit $K$ un hyperrecouvrement de $X$.
Soit $\mathcal{F}$ un faisceau de groupes abéliens sur $\mathcal{C}$.
Il existe un morphisme
$$
s(\mathcal{F}(K))
\longrightarrow
R\Gamma(X, \mathcal{F})
$$
dans $D^{+}(\textit{Ab})$, fonctoriel en $\mathcal{F}$, qui induit
des transformations naturelles
$$
\check{H}^i(K, -) \longrightarrow H^i(X, -)
$$
entre foncteurs $\textit{Ab}(\mathcal{C}) \to \textit{Ab}$. De plus,
il existe une suite spectrale $(E_r, d_r)_{r \geq 0}$ avec
$$
E_2^{p, q} = \check{H}^p(K, \underline{H}^q(\mathcal{F}))
$$
convergeant vers $H^{p + q}(X, \mathcal{F})$.
Cette suite spectrale est fonctorielle en $\mathcal{F}$ et
en l'hyperrecouvrement $K$.
\end{lemma}

\begin{proof}
On pourrait le démontrer par la même méthode que celle employée dans le lemme
correspondant du chapitre consacré à la cohomologie. Donnons plutôt une
démonstration par un complexe double.

\medskip\noindent
Choisissons une résolution injective $\mathcal{F} \to \mathcal{I}^\bullet$
dans la catégorie des faisceaux abéliens sur $\mathcal{C}$. Considérons le
complexe double $A^{\bullet, \bullet}$ de termes
$$
A^{p, q} = \mathcal{I}^q(K_p)
$$
où la différentielle $d_1^{p, q} : A^{p, q} \to A^{p + 1, q}$ est celle
qui provient de la différentielle du complexe $s(\mathcal{I}^q(K))$
associé au groupe abélien cosimplicial $\mathcal{I}^p(K)$,
et où la différentielle $d_2^{p, q} : A^{p, q} \to A^{p, q + 1}$ est celle
qui provient de la différentielle $\mathcal{I}^q \to \mathcal{I}^{q + 1}$.
Notons $\text{Tot}(A^{\bullet, \bullet})$ le complexe total associé au
complexe double $A^{\bullet, \bullet}$ ; voir
Homologie, section \ref{homology-section-double-complexes}.
Nous utiliserons les deux suites
spectrales $({}'E_r, {}'d_r)$ et $({}''E_r, {}''d_r)$
associées à ce complexe double ; voir
Homologie, section \ref{homology-section-double-complex}.

\medskip\noindent
D'après le Lemme \ref{lemma-injective-trivial-cech},
les complexes $s(\mathcal{I}^q(K))$ sont acycliques en
degrés positifs et ont pour $H^0$ le groupe $\mathcal{I}^q(X)$.
Par conséquent, d'après
Homologie, Lemme \ref{homology-lemma-double-complex-gives-resolution},
le morphisme naturel
$$
\mathcal{I}^\bullet(X) \longrightarrow \text{Tot}(A^{\bullet, \bullet})
$$
est un quasi-isomorphisme de complexes de groupes abéliens. En particulier,
on conclut que $H^n(\text{Tot}(A^{\bullet, \bullet})) = H^n(X, \mathcal{F})$.

\medskip\noindent
Le morphisme $s(\mathcal{F}(K)) \longrightarrow R\Gamma(X, \mathcal{F})$ du
lemme est le composé du morphisme
$s(\mathcal{F}(K)) \to \text{Tot}(A^{\bullet, \bullet})$
avec l'inverse
du quasi-isomorphisme affiché ci-dessus. Cela convient puisque
$\mathcal{I}^\bullet(X)$ représente $R\Gamma(X, \mathcal{F})$.

\medskip\noindent
Considérons la suite spectrale $({}'E_r, {}'d_r)_{r \geq 0}$. D'après
Homologie, Lemme \ref{homology-lemma-ss-double-complex},
on voit que
$$
{}'E_2^{p, q} = H^p_I(H^q_{II}(A^{\bullet, \bullet}))
$$
Autrement dit, on calcule d'abord la cohomologie relativement à
$d_2$, ce qui donne les groupes
${}'E_1^{p, q} = \underline{H}^q(\mathcal{F})(K_p)$.
Il en résulte bien (d'après la description de la différentielle
${}'d_1$) que
${}'E_2^{p, q} = \check{H}^p(K, \underline{H}^q(\mathcal{F}))$.
De ce qui précède et de Homologie, Lemme \ref{homology-lemma-first-quadrant-ss},
on déduit que cette suite converge vers $H^n(X, \mathcal{F})$, comme voulu.

\medskip\noindent
Nous omettons la démonstration des assertions relatives à la fonctorialité des
constructions précédentes en le faisceau abélien $\mathcal{F}$ et en
l'hyperrecouvrement $K$.
\end{proof}









\section{Hyperrecouvrements à la Verdier}
\label{section-hypercoverings-verdier}

\noindent
Le lecteur attentif aura remarqué que, pour obtenir
la suite spectrale de {\v C}ech à la cohomologie associée à un
hyperrecouvrement d'un objet $X$, il suffit de disposer de la
conclusion du Lemme \ref{lemma-hypercovering-F}.
La définition suivante a donc un sens.

\begin{definition}
\label{definition-hypercovering-variant}
Soit $\mathcal{C}$ un site. Supposons que $\mathcal{C}$ admette des égalisateurs
et des produits fibrés. Soit $\mathcal{G}$ un préfaisceau d'ensembles.
Un {\it hyperrecouvrement de $\mathcal{G}$} est un objet simplicial
$K$ de $\text{SR}(\mathcal{C})$, muni d'une augmentation
$F(K) \to \mathcal{G}$ telle que
\begin{enumerate}
\item $F(K_0) \to \mathcal{G}$ devienne surjectif
après passage aux faisceaux associés,
\item $F(K_1) \to F(K_0) \times_\mathcal{G} F(K_0)$
devienne surjectif après passage aux faisceaux associés, et
\item $F(K_{n + 1}) \longrightarrow F((\text{cosk}_n \text{sk}_n K)_{n + 1})$
devienne surjectif après passage aux faisceaux associés pour $n \geq 1$.
\end{enumerate}
On dit qu'un objet simplicial $K$ de $\text{SR}(\mathcal{C})$
est un {\it hyperrecouvrement} si $K$ est un hyperrecouvrement de l'objet
final $*$ de $\textit{PSh}(\mathcal{C})$.
\end{definition}

\noindent
L'hypothèse selon laquelle $\mathcal{C}$ admet des produits fibrés et des égalisateurs
garantit que $\text{SR}(\mathcal{C})$ admet des produits fibrés
et des égalisateurs, et que $F$ commute à ceux-ci
(Lemme \ref{lemma-coprod-prod-SR}), ce qui suffit
pour définir les foncteurs cosquelette utilisés (voir
Simplicial, Remarque \ref{simplicial-remark-existence-cosk} et
Catégories, Lemme \ref{categories-lemma-fibre-products-equalizers-exist}).
Si $\mathcal{C}$ est quelconque, on peut remplacer la condition (3) par la
condition selon laquelle
$F(K_{n + 1}) \longrightarrow ((\text{cosk}_n \text{sk}_n F(K))_{n + 1})$
devient surjectif après passage aux faisceaux associés pour $n \geq 1$, et les
résultats de cette section restent valables.

\medskip\noindent
Soit $\mathcal{F}$ un faisceau abélien sur $\mathcal{C}$.
Dans la section précédente, nous avons défini le complexe de {\v C}ech de $\mathcal{F}$
relativement à un objet simplicial $K$ de $\text{SR}(\mathcal{C})$.
Puis, étant donné un préfaisceau $\mathcal{G}$, on pose
$$
H^0(\mathcal{G}, \mathcal{F}) =
\Mor_{\textit{PSh}(\mathcal{C})}(\mathcal{G}, \mathcal{F}) =
\Mor_{\Sh(\mathcal{C})}(\mathcal{G}^\#, \mathcal{F}) =
H^0(\mathcal{G}^\#, \mathcal{F})
$$
avec les notations de
Cohomologie sur les sites, section \ref{sites-cohomology-section-limp}.
C'est un foncteur exact à gauche, et ses foncteurs dérivés supérieurs
(brièvement étudiés dans
Cohomologie sur les sites, section \ref{sites-cohomology-section-limp})
sont notés $H^i(\mathcal{G}, \mathcal{F})$.
Nous montrerons qu'étant donné un hyperrecouvrement $K$ de $\mathcal{G}$,
il existe une suite spectrale de {\v C}ech à la cohomologie convergeant vers la
cohomologie $H^i(\mathcal{G}, \mathcal{F})$.
Remarquons que si $\mathcal{G} = *$, alors
$H^i(*, \mathcal{F}) = H^i(\mathcal{C}, \mathcal{F})$ redonne
la cohomologie de $\mathcal{F}$ sur le site $\mathcal{C}$.

\begin{lemma}
\label{lemma-h0-cech-variant}
Soit $\mathcal{C}$ un site admettant des égalisateurs et des produits fibrés.
Soit $\mathcal{G}$ un préfaisceau sur $\mathcal{C}$.
Soit $K$ un hyperrecouvrement de $\mathcal{G}$.
Soit $\mathcal{F}$ un faisceau de groupes abéliens sur $\mathcal{C}$.
Alors $\check{H}^0(K, \mathcal{F}) = H^0(\mathcal{G}, \mathcal{F})$.
\end{lemma}

\begin{proof}
Cela résulte de la définition de $H^0(\mathcal{G}, \mathcal{F})$
et du fait que
$$
\xymatrix{
F(K_1) \ar@<1ex>[r] \ar@<-1ex>[r] &
F(K_0) \ar[r] & \mathcal{G}
}
$$
devient un diagramme coégalisateur après passage aux faisceaux associés.
\end{proof}

\begin{lemma}
\label{lemma-injective-trivial-cech-variant}
Soit $\mathcal{C}$ un site admettant des égalisateurs et des produits fibrés.
Soit $\mathcal{G}$ un préfaisceau sur $\mathcal{C}$.
Soit $K$ un hyperrecouvrement de $\mathcal{G}$.
Soit $\mathcal{I}$ un faisceau injectif de groupes abéliens sur $\mathcal{C}$.
Alors
$$
\check{H}^p(K, \mathcal{I}) =
\left\{
\begin{matrix}
H^0(\mathcal{G}, \mathcal{I}) & \text{si} & p = 0 \\
0 & \text{si} & p > 0
\end{matrix}
\right.
$$
\end{lemma}

\begin{proof}
D'après (\ref{equation-identify-cech}), on a
$$
s(\mathcal{F}(K))
=
\Hom_{\textit{Ab}(\mathcal{C})}(s(\mathbf{Z}_{F(K)}^\#), \mathcal{F})
$$
Le complexe $s(\mathbf{Z}_{F(K)}^\#)$ est quasi-isomorphe
à $\mathbf{Z}_\mathcal{G}^\#$ ; voir le
Lemme \ref{lemma-acyclic-hypercover-sheaves}. On en conclut
que si $\mathcal{I}$ est un faisceau abélien injectif, alors
le complexe $s(\mathcal{I}(K))$ est acyclique, sauf peut-être en degré $0$.
Autrement dit, on a $\check{H}^i(K, \mathcal{I}) = 0$
pour $i > 0$. Combiné au Lemme \ref{lemma-h0-cech-variant},
ceci démontre le lemme.
\end{proof}

\begin{lemma}
\label{lemma-cech-spectral-sequence-variant}
Soit $\mathcal{C}$ un site admettant des égalisateurs et des produits fibrés.
Soit $\mathcal{G}$ un préfaisceau sur $\mathcal{C}$.
Soit $K$ un hyperrecouvrement de $\mathcal{G}$.
Soit $\mathcal{F}$ un faisceau de groupes abéliens sur $\mathcal{C}$.
Il existe un morphisme
$$
s(\mathcal{F}(K)) \longrightarrow R\Gamma(\mathcal{G}, \mathcal{F})
$$
dans $D^{+}(\textit{Ab})$, fonctoriel en $\mathcal{F}$, qui induit
une transformation naturelle
$$
\check{H}^i(K, -) \longrightarrow H^i(\mathcal{G}, -)
$$
de foncteurs $\textit{Ab}(\mathcal{C}) \to \textit{Ab}$. De plus,
il existe une suite spectrale $(E_r, d_r)_{r \geq 0}$ avec
$$
E_2^{p, q} = \check{H}^p(K, \underline{H}^q(\mathcal{F}))
$$
convergeant vers $H^{p + q}(\mathcal{G}, \mathcal{F})$.
Cette suite spectrale est fonctorielle en $\mathcal{F}$ et
en l'hyperrecouvrement $K$.
\end{lemma}

\begin{proof}
Choisissons une résolution injective $\mathcal{F} \to \mathcal{I}^\bullet$
dans la catégorie des faisceaux abéliens sur $\mathcal{C}$. Considérons le
complexe double $A^{\bullet, \bullet}$ de termes
$$
A^{p, q} = \mathcal{I}^q(K_p)
$$
où la différentielle $d_1^{p, q} : A^{p, q} \to A^{p + 1, q}$
est celle provenant de la différentielle $\mathcal{I}^p \to \mathcal{I}^{p + 1}$,
et où la différentielle $d_2^{p, q} : A^{p, q} \to A^{p, q + 1}$ est
celle provenant de la différentielle du complexe
$s(\mathcal{I}^p(K))$ associé au groupe abélien cosimplicial
$\mathcal{I}^p(K)$, comme expliqué ci-dessus.
Nous utiliserons les deux suites
spectrales $({}'E_r, {}'d_r)$ et $({}''E_r, {}''d_r)$
associées à ce complexe double ; voir
Homologie, section \ref{homology-section-double-complex}.

\medskip\noindent
D'après le Lemme \ref{lemma-injective-trivial-cech-variant}, les complexes
$s(\mathcal{I}^p(K))$ sont acycliques en degrés positifs et ont pour
$H^0$ le groupe $H^0(\mathcal{G}, \mathcal{I}^p)$. Par conséquent, d'après
Homologie, Lemme \ref{homology-lemma-double-complex-gives-resolution},
et sa démonstration, la suite spectrale $({}'E_r, {}'d_r)$ dégénère,
et le morphisme naturel
$$
H^0(\mathcal{G}, \mathcal{I}^\bullet) \longrightarrow
\text{Tot}(A^{\bullet, \bullet})
$$
est un quasi-isomorphisme de complexes de groupes abéliens. Le morphisme
$s(\mathcal{F}(K)) \longrightarrow R\Gamma(\mathcal{G}, \mathcal{F})$
du lemme est le composé du morphisme naturel
$s(\mathcal{F}(K)) \to \text{Tot}(A^{\bullet, \bullet})$
avec l'inverse du quasi-isomorphisme affiché ci-dessus.
Cela convient parce que $H^0(\mathcal{G}, \mathcal{I}^\bullet)$
représente $R\Gamma(\mathcal{G}, \mathcal{F})$.

\medskip\noindent
Considérons la suite spectrale $({}''E_r, {}''d_r)_{r \geq 0}$. D'après
Homologie, Lemme \ref{homology-lemma-ss-double-complex},
on voit que
$$
{}''E_2^{p, q} = H^p_{II}(H^q_I(A^{\bullet, \bullet}))
$$
Autrement dit, on calcule d'abord la cohomologie relativement à
$d_1$, ce qui donne les groupes
${}''E_1^{p, q} = \underline{H}^p(\mathcal{F})(K_q)$.
Il en résulte bien (d'après la description de la différentielle
${}''d_1$) que
${}''E_2^{p, q} = \check{H}^p(K, \underline{H}^q(\mathcal{F}))$.
Puisque cette suite spectrale converge vers la cohomologie de
$\text{Tot}(A^{\bullet, \bullet})$, la démonstration est terminée.
\end{proof}

\begin{lemma}
\label{lemma-cech-spectral-sequence-verdier}
Soit $\mathcal{C}$ un site admettant des égalisateurs et des produits fibrés.
Soit $K$ un hyperrecouvrement.
Soit $\mathcal{F}$ un faisceau abélien. Il existe une
suite spectrale $(E_r, d_r)_{r \geq 0}$ avec
$$
E_2^{p, q} = \check{H}^p(K, \underline{H}^q(\mathcal{F}))
$$
convergeant vers les groupes de cohomologie globale $H^{p + q}(\mathcal{F})$.
\end{lemma}

\begin{proof}
C'est un cas particulier du Lemme \ref{lemma-cech-spectral-sequence-variant}.
\end{proof}







\section{Recouvrement des hyperrecouvrements}
\label{section-covering}

\noindent
Voici quelques méthodes pour construire des hyperrecouvrements.
Remarquons que, puisque la catégorie
$\text{SR}(\mathcal{C}, X)$ admet des produits fibrés,
la catégorie des objets simpliciaux
de $\text{SR}(\mathcal{C}, X)$ en admet elle aussi ;
voir Simplicial, Lemme \ref{simplicial-lemma-fibre-product}.

\begin{lemma}
\label{lemma-funny-gamma}
Soit $\mathcal{C}$ un site admettant des produits fibrés.
Soit $X$ un objet de $\mathcal{C}$.
Soient $K, L, M$ des objets simpliciaux de $\text{SR}(\mathcal{C}, X)$.
Soient $a : K \to L$, $b : M \to L$ des morphismes.
Supposons que
\begin{enumerate}
\item $K$ soit un hyperrecouvrement de $X$,
\item le morphisme $M_0 \to L_0$ soit un recouvrement, et
\item pour tout $n \geq 0$, dans le diagramme
$$
\xymatrix{
M_{n + 1} \ar[dd] \ar[rr] \ar[rd]^\gamma &
&
(\text{cosk}_n \text{sk}_n M)_{n + 1} \ar[dd] \\
&
L_{n + 1}
\times_{(\text{cosk}_n \text{sk}_n L)_{n + 1}}
(\text{cosk}_n \text{sk}_n M)_{n + 1}
\ar[ld] \ar[ru]
& \\
L_{n + 1} \ar[rr] & & (\text{cosk}_n \text{sk}_n L)_{n + 1}
}
$$
la flèche $\gamma$ soit un recouvrement.
\end{enumerate}
Alors le produit fibré $K \times_L M$ est un hyperrecouvrement de $X$.
\end{lemma}

\begin{proof}
Le morphisme $(K \times_L M)_0 = K_0 \times_{L_0} M_0 \to K_0$
est le changement de base d'un recouvrement d'après (2), donc un recouvrement ; voir
le Lemme \ref{lemma-covering-permanence}. Et $K_0 \to \{X \to X\}$
est un recouvrement d'après (1). Ainsi, $(K \times_L M)_0 \to \{X \to X\}$
est un recouvrement d'après le Lemme \ref{lemma-covering-permanence}. Par conséquent,
$K \times_L M$ satisfait la première condition de la Définition
\ref{definition-hypercovering}.

\medskip\noindent
Il reste à vérifier que
$$
K_{n + 1} \times_{L_{n + 1}} M_{n + 1} = (K \times_L M)_{n + 1}
\longrightarrow
(\text{cosk}_n \text{sk}_n (K \times_L M))_{n + 1}
$$
est un recouvrement pour tout $n \geq 0$. Adoptons les abréviations suivantes :
$A = (\text{cosk}_n \text{sk}_n K)_{n + 1}$,
$B = (\text{cosk}_n \text{sk}_n L)_{n + 1}$, et
$C = (\text{cosk}_n \text{sk}_n M)_{n + 1}$.
Le foncteur $\text{cosk}_n \text{sk}_n$ commute aux produits fibrés ;
voir Simplicial, Lemme \ref{simplicial-lemma-cosk-fibre-product}.
Le membre de droite ci-dessus est donc égal à $A \times_B C$.
Considérons le diagramme commutatif suivant
$$
\xymatrix{
K_{n + 1} \times_{L_{n + 1}} M_{n + 1} \ar[r] \ar[d] &
M_{n + 1} \ar[d] \ar[rd]_\gamma \ar[rrd] &
& \\
K_{n + 1} \ar[r] \ar[rd] &
L_{n + 1} \ar[rrd] &
L_{n + 1} \times_B C \ar[l] \ar[r] &
C \ar[d] \\
&
A \ar[rr] &
&
B
}
$$
Ce diagramme montre que
$$
K_{n + 1} \times_{L_{n + 1}} M_{n + 1}
=
(K_{n + 1} \times_B C)
\times_{(L_{n + 1} \times_B C), \gamma}
M_{n + 1}
$$
Or $K_{n + 1} \times_B C \to A \times_B C$
est le changement de base du recouvrement $K_{n + 1} \to A$
par le morphisme $A \times_B C \to A$ ; c'est donc un
recouvrement. D'après l'hypothèse (3), le morphisme $\gamma$ est un recouvrement.
Par conséquent, le morphisme
$$
(K_{n + 1} \times_B C)
\times_{(L_{n + 1} \times_B C), \gamma}
M_{n + 1}
\longrightarrow
K_{n + 1} \times_B C
$$
est un recouvrement, puisqu'il est le changement de base d'un recouvrement.
Le lemme en découle, car un composé de recouvrements
est un recouvrement.
\end{proof}

\begin{lemma}
\label{lemma-product-hypercoverings}
Soit $\mathcal{C}$ un site admettant des produits fibrés.
Soit $X$ un objet de $\mathcal{C}$.
Si $K, L$ sont des hyperrecouvrements de $X$, alors
$K \times L$ est un hyperrecouvrement de $X$.
\end{lemma}

\begin{proof}
On peut soit le vérifier directement, soit appliquer le
Lemme \ref{lemma-funny-gamma} ci-dessus et vérifier que $L \to \{X \to X\}$
possède la propriété (3).
\end{proof}


\noindent
Soit $\mathcal{C}$ un site admettant des produits fibrés.
Soit $X$ un objet de $\mathcal{C}$.
Puisque la catégorie $\text{SR}(\mathcal{C}, X)$ admet des coproduits et des
limites finies, il est permis de parler des objets
$U \times K$ et $\Hom(U, K)$ pour certains ensembles simpliciaux $U$
(par exemple ceux qui ne possèdent qu'un nombre fini de simplexes non dégénérés)
et tout objet simplicial $K$ de $\text{SR}(\mathcal{C}, X)$.
Voir Simplicial, sections
\ref{simplicial-section-product-with-simplicial-sets} et
\ref{simplicial-section-hom-from-simplicial-sets}.

\begin{lemma}
\label{lemma-covering}
Soit $\mathcal{C}$ un site admettant des produits fibrés.
Soit $X$ un objet de $\mathcal{C}$.
Soit $K$ un hyperrecouvrement de $X$.
Soit $k \geq 0$ un entier.
Soit $u : Z \to K_k$ un recouvrement
dans $\text{SR}(\mathcal{C}, X)$.
Alors il existe un morphisme d'hyperrecouvrements
$f: L \to K$ tel que $L_k \to K_k$
se factorise par $u$.
\end{lemma}

\begin{proof}
Notons $Y = K_k$. Soit $C[k]$ l'ensemble cosimplicial défini dans
Simplicial, Exemple \ref{simplicial-example-simplex-cosimplicial-set}.
Nous utiliserons la description de $\Hom(C[k], Y)$ et $\Hom(C[k], Z)$
donnée dans
Simplicial, Lemme \ref{simplicial-lemma-morphism-into-product}.
Il existe un morphisme canonique
$K \to \Hom(C[k], Y)$ correspondant à $\text{id} : K_k = Y \to Y$.
Considérons le morphisme $\Hom(C[k], Z) \to \Hom(C[k], Y)$
qui, sur les termes de degré $n$, est le morphisme
$$
\prod\nolimits_{\alpha : [k] \to [n]} Z
\longrightarrow
\prod\nolimits_{\alpha : [k] \to [n]} Y
$$
obtenu en utilisant le morphisme donné $Z \to Y$ sur chaque facteur. Posons
$$
L = K \times_{\Hom(C[k], Y)} \Hom(C[k], Z).
$$
Le morphisme $L_k \to K_k$ s'insère dans un diagramme commutatif
$$
\xymatrix{
L_k \ar[r] \ar[d] &
\prod_{\alpha : [k] \to [k]} Z \ar[r]^-{\text{pr}_{\text{id}_{[k]}}} \ar[d] &
Z \ar[d] \\
K_k \ar[r] &
\prod_{\alpha : [k] \to [k]} Y \ar[r]^-{\text{pr}_{\text{id}_{[k]}}} &
Y
}
$$
Comme le composé des deux flèches du bas est l'identité,
on en conclut que l'on a la factorisation voulue.

\medskip\noindent
Il reste à montrer que $L$ est un hyperrecouvrement de $X$.
Pour cela, nous appliquerons le Lemme \ref{lemma-funny-gamma}.
La condition (1) est satisfaite par hypothèse.
Pour (2), le morphisme
$$
\Hom(C[k], Z)_0 \to \Hom(C[k], Y)_0
$$
est un recouvrement parce qu'il est isomorphe à $Z \to Y$, puisqu'il
n'existe qu'un seul morphisme $[k] \to [0]$.

\medskip\noindent
Considérons la condition (3) pour $n = 0$. Alors, puisque
$(\text{cosk}_0 T)_1 = T \times T$
(Simplicial, Exemple \ref{simplicial-example-cosk0})
et puisque $\Hom(C[k], Z)_1 = \prod_{\alpha : [k] \to [1]} Z$,
on obtient le diagramme
$$
\xymatrix{
\prod\nolimits_{\alpha : [k] \to [1]} Z \ar[r] \ar[d] &
Z \times Z \ar[d] \\
\prod\nolimits_{\alpha : [k] \to [1]} Y \ar[r] &
Y \times Y
}
$$
où les flèches horizontales correspondent à la projection sur les facteurs
associés aux deux $\alpha$ non surjectifs. La flèche $\gamma$
est donc le morphisme
$$
\prod\nolimits_{\alpha : [k] \to [1]} Z
\longrightarrow
\prod\nolimits_{\alpha : [k] \to [1]\text{ non surjective}} Z
\times
\prod\nolimits_{\alpha : [k] \to [1]\text{ surjective}} Y
$$
qui est un produit de recouvrements, donc un recouvrement d'après le
Lemme \ref{lemma-covering-permanence}.

\medskip\noindent
Considérons la condition (3) pour $n > 0$. Nous affirmons qu'il existe une
application injective $\tau : S' \to S$ d'ensembles finis telle que, pour tout
objet $T$ de $\text{SR}(\mathcal{C}, X)$, le morphisme
\begin{equation}
\label{equation-map}
\Hom(C[k], T)_{n + 1} \to
(\text{cosk}_n \text{sk}_n \Hom(C[k], T))_{n + 1}
\end{equation}
soit isomorphe à la projection $\prod_{s \in S} T \to \prod_{s' \in S'} T$,
fonctoriellement en $T$. Si cela est vrai, on voit, en raisonnant comme au
paragraphe précédent, que la flèche $\gamma$ est le morphisme
$$
\prod\nolimits_{s \in S} Z
\longrightarrow
\prod\nolimits_{s \in S'} Z
\times
\prod\nolimits_{s \not\in \tau(S')} Y
$$
qui est un produit de recouvrements, donc un recouvrement d'après le
Lemme \ref{lemma-covering-permanence}. Par construction, on a
$\Hom(C[k], T)_{n + 1} = \prod_{\alpha : [k] \to [n + 1]} T$
(voir Simplicial, Lemme \ref{simplicial-lemma-morphism-into-product}).
On prend donc $S = \text{Map}([k], [n + 1])$.
D'autre part, Simplicial, Lemme \ref{simplicial-lemma-formula-limit},
donne une description des points de
$(\text{cosk}_n \text{sk}_n \Hom(C[k], T))_{n + 1}$
comme des suites $(f_0, \ldots, f_{n + 1})$ de points de $\Hom(C[k], T)_n$
vérifiant $d^n_{j - 1} f_i = d^n_i f_j$ pour $0 \leq i < j \leq n + 1$.
On peut écrire $f_i = (f_{i, \alpha})$, où $f_{i, \alpha}$ est un point de $T$
et $\alpha \in \text{Map}([k], [n])$. Les conditions se traduisent par
$$
f_{i, \delta^n_{j - 1} \circ \beta} = f_{j, \delta_i^n \circ \beta}
$$
pour tout $0 \leq i < j \leq n + 1$ et tout $\beta : [k] \to [n - 1]$. On
voit donc que
$$
S' = \{0, \ldots, n + 1\} \times \text{Map}([k], [n]) / \sim
$$
où la relation d'équivalence est engendrée par les équivalences
$$
(i, \delta^n_{j - 1} \circ \beta) \sim (j, \delta_i^n \circ \beta)
$$
pour $0 \leq i < j \leq n + 1$ et $\beta : [k] \to [n - 1]$.
Un calcul (omis) montre que le morphisme (\ref{equation-map})
correspond à l'application $S' \to S$ qui envoie $(i, \alpha)$ sur
$\delta^{n + 1}_i \circ \alpha \in S$. (Le lecteur pourra être rassuré
de voir que cette application est bien définie d'après la partie (1) du
Lemme \ref{simplicial-lemma-relations-face-degeneracy} de Simplicial.)
Pour achever la démonstration, il suffit de montrer que si
$\alpha, \alpha' : [k] \to [n]$ et $0 \leq i < j \leq n + 1$
sont tels que
$$
\delta^{n + 1}_i \circ \alpha = \delta^{n + 1}_j \circ \alpha'
$$
alors on a $\alpha = \delta^n_{j - 1} \circ \beta$
et $\alpha' = \delta_i^n \circ \beta$ pour un certain $\beta : [k] \to [n - 1]$.
Cela se voit aisément, et nous l'omettons.
\end{proof}

\begin{lemma}
\label{lemma-covering-sheaf}
Soit $\mathcal{C}$ un site admettant des produits fibrés.
Soit $X$ un objet de $\mathcal{C}$.
Soit $K$ un hyperrecouvrement de $X$.
Soit $n \geq 0$ un entier.
Soit $u : \mathcal{F} \to F(K_n)$ un morphisme
de préfaisceaux qui devient surjectif
après passage aux faisceaux associés.
Alors il existe un morphisme d'hyperrecouvrements
$f: L \to K$ tel que $F(f_n) : F(L_n) \to F(K_n)$
se factorise par $u$.
\end{lemma}

\begin{proof}
Écrivons $K_n = \{U_i \to X\}_{i \in I}$.
Ainsi, l'application $u$ est un morphisme de préfaisceaux d'ensembles
$u : \mathcal{F} \to \amalg h_{u_i}$.
L'hypothèse sur $u$ signifie que, pour tout
$i \in I$, il existe un recouvrement $\{U_{ij} \to U_i\}_{j \in I_i}$
du site $\mathcal{C}$ et un morphisme de préfaisceaux
$t_{ij} : h_{U_{ij}} \to \mathcal{F}$ tel que
$u \circ t_{ij}$ soit l'application $h_{U_{ij}} \to h_{U_i}$
provenant du morphisme $U_{ij} \to U_i$.
Posons $J = \amalg_{i \in I} I_i$, et soit
$\alpha : J \to I$ l'application évidente.
Pour $j \in J$, notons $V_j = U_{\alpha(j)j}$. Posons
$Z = \{V_j \to X\}_{j \in J}$.
Enfin, considérons le morphisme
$u' : Z \to K_n$ donné par $\alpha : J \to I$
et par les morphismes $V_j = U_{\alpha(j)j} \to U_{\alpha(j)}$
ci-dessus. Il est clair que c'est un recouvrement dans la
catégorie $\text{SR}(\mathcal{C}, X)$ et, par
construction, $F(u') : F(Z) \to F(K_n)$ se factorise par $u$.
Le résultat découle donc du Lemme \ref{lemma-covering} ci-dessus.
\end{proof}


\section{Ajout de simplexes}
\label{section-adding-simplices}

\noindent
Dans cette section, nous démontrons quelques lemmes techniques qui serviront plus loin.
Soit $\mathcal{C}$ un site admettant des produits fibrés.
Soit $X$ un objet de $\mathcal{C}$.
Comme nous l'avons signalé dans la section \ref{section-covering} ci-dessus,
les objets $U \times K$ et $\Hom(U, K)$
sont définis pour certains ensembles simpliciaux $U$
et tout objet simplicial $K$ de $\text{SR}(\mathcal{C}, X)$.
Voir Simplicial, sections
\ref{simplicial-section-product-with-simplicial-sets} et
\ref{simplicial-section-hom-from-simplicial-sets}.

\begin{lemma}
\label{lemma-one-more-simplex}
Soit $\mathcal{C}$ un site admettant des produits fibrés.
Soit $X$ un objet de $\mathcal{C}$.
Soit $K$ un hyperrecouvrement de $X$.
Soient $U \subset V$ des ensembles simpliciaux tels que $U_n, V_n$
soient finis non vides pour tout $n$.
Supposons que $U$ n'ait qu'un nombre fini de simplexes non dégénérés.
Supposons que $n \geq 0$ et que $x \in V_n$,
$x \not \in U_n$, soient tels que
\begin{enumerate}
\item $V_i = U_i$ pour $i < n$,
\item $V_n = U_n \cup \{x\}$,
\item tout $z \in V_j$, $z \not \in U_j$, pour $j > n$,
soit dégénéré.
\end{enumerate}
Alors le morphisme
$$
\Hom(V, K)_0
\longrightarrow
\Hom(U, K)_0
$$
de $\text{SR}(\mathcal{C}, X)$ est un recouvrement.
\end{lemma}

\begin{proof}
Si $n = 0$, il résulte aisément que $V = U \amalg \Delta[0]$
(voir ci-dessous). Dans ce cas, $\Hom(V, K)_0 =
\Hom(U, K)_0 \times K_0$. Le résultat découle alors
du Lemme \ref{lemma-covering-permanence}.

\medskip\noindent
Soit $a : \Delta[n] \to V$ le morphisme associé à $x$
comme dans Simplicial, Lemme \ref{simplicial-lemma-simplex-map}.
Écrivons $\partial \Delta[n] = i_{(n-1)!} \text{sk}_{n - 1} \Delta[n]$
pour le $(n - 1)$-squelette de $\Delta[n]$.
Soit $b : \partial \Delta[n] \to U$ la restriction
de $a$ au $(n - 1)$-squelette de $\Delta[n]$. D'après
Simplicial, Lemme \ref{simplicial-lemma-glue-simplex},
on a $V = U \amalg_{\partial \Delta[n]} \Delta[n]$. D'après
Simplicial, Lemme
\ref{simplicial-lemma-hom-from-coprod},
le carré
$$
\xymatrix{
\Hom(V, K)_0 \ar[r] \ar[d] &
\Hom(U, K)_0 \ar[d] \\
\Hom(\Delta[n], K)_0 \ar[r] &
\Hom(\partial \Delta[n], K)_0
}
$$
est cartésien. Il suffit donc de montrer que
la flèche horizontale inférieure est un recouvrement. D'après
Simplicial, Lemme \ref{simplicial-lemma-cosk-shriek},
cette flèche s'identifie à
$$
K_n \to (\text{cosk}_{n - 1} \text{sk}_{n - 1} K)_n
$$
et c'est donc un recouvrement par définition d'un hyperrecouvrement.
\end{proof}

\begin{lemma}
\label{lemma-add-simplices}
Soit $\mathcal{C}$ un site admettant des produits fibrés.
Soit $X$ un objet de $\mathcal{C}$.
Soit $K$ un hyperrecouvrement de $X$.
Soient $U \subset V$ des ensembles simpliciaux tels que $U_n, V_n$
soient finis non vides pour tout $n$.
Supposons que $U$ et $V$ n'aient qu'un nombre fini de simplexes non dégénérés.
Alors le morphisme
$$
\Hom(V, K)_0
\longrightarrow
\Hom(U, K)_0
$$
de $\text{SR}(\mathcal{C}, X)$ est un recouvrement.
\end{lemma}

\begin{proof}
D'après le Lemme \ref{lemma-one-more-simplex}
ci-dessus, il suffit de démontrer un lemme élémentaire
sur les inclusions d'ensembles simpliciaux $U \subset V$ comme dans le
lemme. Or c'est exactement le résultat de
Simplicial, Lemme \ref{simplicial-lemma-add-simplices}.
\end{proof}

\begin{lemma}
\label{lemma-degeneracy-maps-coverings}
Soit $\mathcal{C}$ un site admettant des produits fibrés. Soit $X$ un objet de
$\mathcal{C}$. Soit $K$ un hyperrecouvrement de $X$. Alors
\begin{enumerate}
\item $K_n$ est un recouvrement de $X$ pour tout $n \geq 0$,
\item $d^n_i : K_n \to K_{n - 1}$ est un recouvrement pour tous $n \geq 1$
et $0 \leq i \leq n$.
\end{enumerate}
\end{lemma}

\begin{proof}
Rappelons que $K_0$ est un recouvrement de $X$ d'après la
Définition \ref{definition-hypercovering}
et que cela équivaut à dire que
$K_0 \to \{X \to X\}$ est un recouvrement au sens
de la Définition \ref{definition-covering-SR}.
L'assertion (1) résulte donc de (2), car celle-ci montrera que
le composé
$K_n \to K_{n - 1} \to \ldots \to K_0 \to \{X \to X\}$
est un recouvrement d'après le Lemme \ref{lemma-covering-permanence}.

\medskip\noindent
Démonstration de (2). Observons que
$\Mor(\Delta[n], K)_0 = K_n$ d'après
Simplicial, Lemme \ref{simplicial-lemma-exists-hom-from-simplicial-set-finite}.
L'assertion (2) résulte donc du Lemme \ref{lemma-add-simplices}
appliqué aux $n + 1$ inclusions distinctes $\Delta[n - 1] \to \Delta[n]$.
\end{proof}

\begin{remark}
\label{remark-P-covering}
Un cas particulier utile des Lemmes \ref{lemma-add-simplices} et
\ref{lemma-degeneracy-maps-coverings} est le suivant.
Supposons donnée une catégorie $\mathcal{C}$ admettant des produits fibrés.
Soit $P \subset \text{Arrows}(\mathcal{C})$ un sous-ensemble
stable par changement de base, stable par composition
et contenant tous les isomorphismes. On appelle alors
{\it $P$-hyperrecouvrement} une augmentation $a : U \to X$
issue d'un objet simplicial de $\mathcal{C}$ telle que
\begin{enumerate}
\item $U_0 \to X$ appartienne à $P$,
\item $U_1 \to U_0 \times_X U_0$ appartienne à $P$,
\item $U_{n + 1} \to (\text{cosk}_n\text{sk}_n U)_{n + 1}$
appartienne à $P$ pour $n \geq 1$.
\end{enumerate}
La catégorie $\mathcal{C}/X$ admet toutes les limites finies ; les
cosquelettes employés dans la formulation ci-dessus existent donc
(voir Catégories, Lemme \ref{categories-lemma-finite-limits-exist}).
Nous affirmons alors que les morphismes $U_n \to X$ et $d^n_i : U_n \to U_{n - 1}$
appartiennent à $P$. Cela résulte des lemmes précités
en munissant $\mathcal{C}$ de la structure de site dont les recouvrements
sont les $\{f : V \to U\}$ avec $f \in P$, et en prenant pour $K$ l'objet défini par
$K_n = \{U_n \to X\}$.
\end{remark}


\section{Homotopies}
\label{section-homotopies}

\noindent
Soit $\mathcal{C}$ un site admettant des produits fibrés.
Soit $X$ un objet de $\mathcal{C}$.
Soit $L$ un objet simplicial de $\text{SR}(\mathcal{C}, X)$.
D'après
Simplicial, Lemme \ref{simplicial-lemma-exists-hom-from-simplicial-set-finite},
il existe un objet $\Hom(\Delta[1], L)$
de la catégorie $\text{Simp}(\text{SR}(\mathcal{C}, X))$ qui représente le
foncteur
$$
T
\longmapsto
\Mor_{\text{Simp}(\text{SR}(\mathcal{C}, X))}(\Delta[1] \times T, L)
$$
Il existe un morphisme canonique
$$
\Hom(\Delta[1], L) \to L \times L
$$
provenant de $e_i : \Delta[0] \to \Delta[1]$ et de l'identification
$\Hom(\Delta[0], L) = L$.

\begin{lemma}
\label{lemma-hom-hypercovering}
Soit $\mathcal{C}$ un site admettant des produits fibrés.
Soit $X$ un objet de $\mathcal{C}$.
Soit $L$ un objet simplicial de $\text{SR}(\mathcal{C}, X)$.
Soit $n \geq 0$. Considérons le diagramme commutatif
\begin{equation}
\label{equation-diagram}
\xymatrix{
\Hom(\Delta[1], L)_{n + 1} \ar[r] \ar[d] &
(\text{cosk}_n \text{sk}_n \Hom(\Delta[1], L))_{n + 1} \ar[d] \\
(L \times L)_{n + 1} \ar[r] &
(\text{cosk}_n \text{sk}_n (L \times L))_{n + 1}
}
\end{equation}
provenant du morphisme défini ci-dessus.
On peut identifier les termes de ce diagramme comme suit,
où
$\partial \Delta[n + 1] = i_{n!}\text{sk}_n \Delta[n + 1]$
est le $n$-squelette du $(n + 1)$-simplexe :
\begin{eqnarray*}
\Hom(\Delta[1], L)_{n + 1}
& = &
\Hom(\Delta[1] \times \Delta[n + 1], L)_0 \\
(\text{cosk}_n \text{sk}_n \Hom(\Delta[1], L))_{n + 1}
& = &
\Hom(\Delta[1] \times \partial \Delta[n + 1], L)_0 \\
(L \times L)_{n + 1}
& = &
\Hom(
(\Delta[n + 1] \amalg \Delta[n + 1], L)_0 \\
(\text{cosk}_n \text{sk}_n (L \times L))_{n + 1}
& = &
\Hom(
\partial \Delta[n + 1]
\amalg
\partial \Delta[n + 1], L)_0
\end{eqnarray*}
et les morphismes entre ces objets de $\text{SR}(\mathcal{C}, X)$
proviennent du diagramme commutatif d'ensembles simpliciaux
\begin{equation}
\label{equation-dual-diagram}
\xymatrix{
\Delta[1] \times \Delta[n + 1] &
\Delta[1] \times \partial\Delta[n + 1] \ar[l] \\
\Delta[n + 1] \amalg \Delta[n + 1] \ar[u] &
\partial\Delta[n + 1] \amalg \partial\Delta[n + 1]
\ar[l] \ar[u]
}
\end{equation}
De plus, le produit fibré de la flèche inférieure et de la
flèche de droite dans (\ref{equation-diagram}) est égal à
$$
\Hom(U, L)_0
$$
où $U \subset \Delta[1] \times \Delta[n + 1]$
est le plus petit sous-ensemble simplicial tel que
$\Delta[n + 1] \amalg \Delta[n + 1]$ et
$\Delta[1] \times \partial\Delta[n + 1]$ s'y envoient tous deux.
\end{lemma}

\begin{proof}
La première et la troisième égalité sont données par
Simplicial, Lemme \ref{simplicial-lemma-exists-hom-from-simplicial-set-finite}.
La deuxième et la quatrième résultent du lemme cité, combiné au
Lemme \ref{simplicial-lemma-cosk-shriek} de Simplicial.
La dernière assertion résulte du fait que
$U$ est la somme amalgamée de la flèche inférieure et de la flèche droite du
diagramme (\ref{equation-dual-diagram}), d'après
Simplicial, Lemme \ref{simplicial-lemma-hom-from-coprod}.
Pour voir que $U$ est égal à cette somme amalgamée, il suffit
de voir que l'intersection de
$\Delta[n + 1] \amalg \Delta[n + 1]$ et de
$\Delta[1] \times \partial\Delta[n + 1]$
dans $\Delta[1] \times \Delta[n + 1]$ est égale à
$\partial\Delta[n + 1] \amalg \partial\Delta[n + 1]$.
Nous laissons cette vérification au lecteur.
\end{proof}

\begin{lemma}
\label{lemma-homotopy}
Soit $\mathcal{C}$ un site admettant des produits fibrés.
Soit $X$ un objet de $\mathcal{C}$.
Soient $K, L$ des hyperrecouvrements de $X$.
Soient $a, b : K \to L$ des morphismes d'hyperrecouvrements.
Il existe un morphisme d'hyperrecouvrements
$c : K' \to K$ tel que $a \circ c$ soit homotope
à $b \circ c$.
\end{lemma}

\begin{proof}
Considérons le diagramme commutatif suivant
$$
\xymatrix{
K' \ar@{=}[r]^-{def} \ar[rd]_c &
K \times_{(L \times L)} \Hom(\Delta[1], L)
\ar[r] \ar[d] & \Hom(\Delta[1], L) \ar[d] \\
& K \ar[r]^{(a, b)} & L \times L
}
$$
Par la propriété fonctorielle de $\Hom(\Delta[1], L)$,
le composé des morphismes horizontaux
correspond à un morphisme $K' \times \Delta[1] \to L$ qui
définit une homotopie entre $c \circ a$ et $c \circ b$.
Ainsi, si l'on peut montrer que $K'$ est un
hyperrecouvrement de $X$, on obtient le lemme.
Pour cela, nous appliquerons le Lemme \ref{lemma-funny-gamma}
à la paire de morphismes $K \to L \times L$
et $\Hom(\Delta[1], L) \to L \times L$.
La condition (1) du Lemme \ref{lemma-funny-gamma} est satisfaite.
La condition (2) du Lemme \ref{lemma-funny-gamma} est vérifiée parce que
$\Hom(\Delta[1], L)_0 = L_1$, et que le morphisme
$(d^1_0, d^1_1) : L_1 \to L_0 \times L_0$ est un
recouvrement de $\text{SR}(\mathcal{C}, X)$, puisque
$L$ est par hypothèse un hyperrecouvrement.
Pour démontrer la condition (3) du Lemme \ref{lemma-funny-gamma},
on utilise le Lemme \ref{lemma-hom-hypercovering} ci-dessus. D'après
ce lemme, le morphisme $\gamma$ de la condition (3) du Lemme
\ref{lemma-funny-gamma} est le morphisme
$$
\Hom(\Delta[1] \times \Delta[n + 1], L)_0
\longrightarrow
\Hom(U, L)_0
$$
où $U \subset \Delta[1] \times \Delta[n + 1]$.
D'après le Lemme \ref{lemma-add-simplices},
c'est un recouvrement, ce qui démontre l'assertion.
\end{proof}

\begin{remark}
\label{remark-contractible-category}
Remarquons que le point essentiel de la démonstration est l'emploi du
Lemme \ref{lemma-add-simplices}. Ce lemme
est tout à fait général et ne dépend pas de la
forme exacte des ensembles simpliciaux (pourvu qu'ils
n'aient qu'un nombre fini de simplexes non dégénérés).
Il paraît tout à fait raisonnable d'attendre un résultat
du type suivant :
Étant donné un morphisme $a : K \times \partial \Delta[k]
\to L$, où $K$ et $L$ sont des hyperrecouvrements, il
existe un morphisme d'hyperrecouvrements $c : K' \to K$
et un morphisme $g : K' \times \Delta[k] \to L$
tels que
$g|_{K' \times \partial \Delta[k]} =
a \circ (c \times \text{id}_{\partial \Delta[k]})$.
Autrement dit, la catégorie des hyperrecouvrements est
contractile en un sens convenable.
\end{remark}

















\section{Cohomologie et hyperrecouvrements}
\label{section-cohomology}

\noindent
Soit $\mathcal{C}$ un site admettant des produits fibrés.
Soit $X$ un objet de $\mathcal{C}$.
Soit $\mathcal{F}$ un faisceau de groupes abéliens sur $\mathcal{C}$.
Soient $K, L$ des hyperrecouvrements de $X$.
Si $a, b : K \to L$ sont des applications homotopes,
alors $\mathcal{F}(a), \mathcal{F}(b) : \mathcal{F}(K) \to \mathcal{F}(L)$
sont des applications homotopes ; voir
Simplicial, Lemme \ref{simplicial-lemma-functorial-homotopy}.
Elles ont donc le même effet sur les groupes de cohomologie des complexes
de cochaînes associés ; voir
Simplicial, Lemme \ref{simplicial-lemma-homotopy-s-Q}.
Nous allons nous en servir pour définir la limite inductive sur tous les
hyperrecouvrements.

\medskip\noindent
Notons temporairement $\text{HC}(\mathcal{C}, X)$
la catégorie dont les objets sont les hyperrecouvrements de $X$
et dont les morphismes sont les applications entre hyperrecouvrements de $X$
à homotopie près. Nous avons vu qu'il s'agit
d'une catégorie et non d'une « grande » catégorie ;
voir le Lemme \ref{lemma-hypercoverings-set}.
L'opposée de $\text{HC}(\mathcal{C}, X)$
sera la catégorie d'indices de notre diagramme ; voir
Catégories, section \ref{categories-section-limits}, pour la terminologie.
Considérons le diagramme
$$
\check{H}^i(-, \mathcal{F}) :
\text{HC}(\mathcal{C}, X)^{opp}
\longrightarrow
\textit{Ab}.
$$
D'après les Lemmes \ref{lemma-product-hypercoverings} et \ref{lemma-homotopy},
ainsi que la remarque ci-dessus sur les homotopies, ce diagramme est filtrant ; voir
Catégories, Définition \ref{categories-definition-directed}.
La limite inductive
$$
\check{H}^i_{\text{HC}}(X, \mathcal{F}) =
\colim_{K \in \text{HC}(\mathcal{C}, X)} \check{H}^i(K, \mathcal{F})
$$
admet donc une description particulièrement simple (voir la référence citée).

\begin{theorem}
\label{theorem-cohomology-hypercoverings}
Soit $\mathcal{C}$ un site admettant des produits fibrés.
Soit $X$ un objet de $\mathcal{C}$. Soit $i \geq 0$.
Les foncteurs
\begin{eqnarray*}
\textit{Ab}(\mathcal{C}) & \longrightarrow & \textit{Ab} \\
\mathcal{F} & \longmapsto & H^i(X, \mathcal{F}) \\
\mathcal{F} & \longmapsto & \check{H}^i_{\text{HC}}(X, \mathcal{F})
\end{eqnarray*}
sont canoniquement isomorphes.
\end{theorem}

\begin{proof}[Démonstration à l'aide des suites spectrales.]
Supposons que $\xi \in H^p(X, \mathcal{F})$ pour un certain $p \geq 0$.
Montrons que $\xi$ appartient à l'image de l'application
$\check{H}^p(X, \mathcal{F}) \to H^p(X, \mathcal{F})$ du
Lemme \ref{lemma-cech-spectral-sequence}
pour un certain hyperrecouvrement $K$ de $X$.

\medskip\noindent
C'est vrai si $p = 0$ d'après le Lemme \ref{lemma-h0-cech}.
Si $p = 1$, choisissons un hyperrecouvrement de {\v C}ech $K$ de $X$ comme dans
l'Exemple \ref{example-cech}, en partant d'un recouvrement
$K_0 = \{U_i \to X\}$ du site $\mathcal{C}$ tel que
$\xi|_{U_i} = 0$ ; voir
Cohomologie sur les sites,
Lemme \ref{sites-cohomology-lemma-kill-cohomology-class-on-covering}.
Il résulte immédiatement de la suite spectrale
du Lemme \ref{lemma-cech-spectral-sequence} que, dans ce cas, $\xi$ provient
d'un élément de $\check{H}^1(K, \mathcal{F})$.
Dans le cas général, choisissons un hyperrecouvrement quelconque $K$ de $X$
tel que l'image de $\xi$ soit nulle dans $\underline{H}^p(\mathcal{F})(K_0)$
(en utilisant l'Exemple \ref{example-cech} et
Cohomologie sur les sites,
Lemme \ref{sites-cohomology-lemma-kill-cohomology-class-on-covering}
une fois encore).
D'après la suite spectrale du Lemme \ref{lemma-cech-spectral-sequence},
l'obstruction à ce que $\xi$ provienne d'un élément de
$\check{H}^p(K, \mathcal{F})$ est une suite d'éléments
$\xi_1, \ldots, \xi_{p - 1}$ avec
$\xi_q \in \check{H}^{p - q}(K, \underline{H}^q(\mathcal{F}))$
(plus précisément, les images des $\xi_q$ dans certains sous-quotients
de ces groupes).

\medskip\noindent
On peut remplacer par récurrence l'hyperrecouvrement $K$ par des raffinements
de telle sorte que les restrictions des obstructions $\xi_1, \ldots, \xi_{p - 1}$ soient nulles
(et pas seulement leurs images
dans les sous-quotients : il n'y a donc ici aucune subtilité). En effet, supposons être
déjà parvenus à une situation où
$\xi_{q + 1}, \ldots, \xi_{p - 1}$ sont nuls.
Remarquons que $\xi_q \in \check{H}^{p - q}(K, \underline{H}^q(\mathcal{F}))$
est la classe d'un élément
$$
\tilde \xi_q \in
\underline{H}^q(\mathcal{F})(K_{p - q}) =
\prod H^q(U_i, \mathcal{F})
$$
si $K_{p - q} = \{U_i \to X\}_{i \in I}$. Soit $\xi_{q, i}$
la composante de $\tilde \xi_q$ dans $H^q(U_i, \mathcal{F})$.
Comme $q \geq 1$, on peut appliquer
Cohomologie sur les sites,
Lemme \ref{sites-cohomology-lemma-kill-cohomology-class-on-covering},
une fois encore, afin de choisir des recouvrements $\{U_{i, j} \to U_i\}$
du site tels que chaque restriction $\xi_{q, i}|_{U_{i, j}} = 0$.
Considérons l'objet $Z = \{U_{i, j} \to X\}$ de la catégorie
$\text{SR}(\mathcal{C}, X)$ et son morphisme évident
$u : Z \to K_{p - q}$. Il est clair que $u$ est un recouvrement ; voir
la Définition \ref{definition-covering-SR}. D'après le
Lemme \ref{lemma-covering}, il
existe un morphisme $L \to K$ d'hyperrecouvrements de $X$ tel que
$L_{p - q} \to K_{p - q}$ se factorise par $u$. Alors l'image de
$\xi_q$ dans $\underline{H}^q(\mathcal{F})(L_{p - q})$
est nulle. Puisque la suite spectrale du
Lemme \ref{lemma-cech-spectral-sequence}
est fonctorielle, cela signifie qu'après avoir remplacé $K$ par $L$, on atteint la
situation où $\xi_q, \ldots, \xi_{p - 1}$ sont tous nuls.
En poursuivant ainsi, on obtient un hyperrecouvrement où ils sont tous
nuls, et donc $\xi$ appartient à l'image de l'application
$\check{H}^p(X, \mathcal{F}) \to H^p(X, \mathcal{F})$.

\medskip\noindent
Supposons que $K$ soit un hyperrecouvrement de $X$, que
$\xi \in \check{H}^p(K, \mathcal{F})$ et que l'image de
$\xi$ par l'application
$\check{H}^p(X, \mathcal{F}) \to H^p(X, \mathcal{F})$ du
Lemme \ref{lemma-cech-spectral-sequence}
soit nulle. Pour achever la démonstration du théorème, il faut montrer qu'il
existe un morphisme d'hyperrecouvrements $L \to K$ tel que
la restriction de $\xi$ soit nulle dans $\check{H}^p(L, \mathcal{F})$.
D'après la suite spectrale du Lemme \ref{lemma-cech-spectral-sequence},
l'annulation de l'image de $\xi$ dans $H^p(X, \mathcal{F})$
signifie qu'il existe des éléments $\xi_1, \ldots, \xi_{p - 2}$
avec $\xi_q \in \check{H}^{p - 1 - q}(K, \underline{H}^q(\mathcal{F}))$
(plus précisément, leurs images dans certains sous-quotients)
tels que les images $d_{q + 1}^{p - 1 - q, q}\xi_q$ (dans la suite
spectrale) aient pour somme $\xi$. Par le même procédé que ci-dessus,
on peut donc trouver un morphisme d'hyperrecouvrements $L \to K$ tel que
les restrictions des éléments $\xi_q$, $q = 1, \ldots, p - 2$,
dans $\check{H}^{p - 1 - q}(L, \underline{H}^q(\mathcal{F}))$ soient nulles.
Il en résulte que $\xi$ est nul, puisque le morphisme $L \to K$
induit un morphisme de suites spectrales d'après le
Lemme \ref{lemma-cech-spectral-sequence}.
\end{proof}

\begin{proof}[Démonstration sans suites spectrales.]
Nous avons vu le résultat pour $i = 0$ ; voir le Lemme \ref{lemma-h0-cech}.
Nous savons que les foncteurs $H^i(X, -)$ forment un $\delta$-foncteur universel ; voir
Catégories dérivées, Lemme \ref{derived-lemma-higher-derived-functors}.
Pour démontrer le théorème, il suffit de montrer que
la suite de foncteurs $\check{H}^i_{HC}(X, -)$ forme un
$\delta$-foncteur. En effet, nous savons que la cohomologie de {\v C}ech
est nulle sur les faisceaux injectifs (Lemme \ref{lemma-injective-trivial-cech}),
et on peut alors appliquer
Homologie, Lemme \ref{homology-lemma-efface-implies-universal}.

\medskip\noindent
Soit
$$
0 \to \mathcal{F} \to \mathcal{G} \to \mathcal{H} \to 0
$$
une suite exacte courte de faisceaux abéliens sur $\mathcal{C}$.
Soit $\xi \in \check{H}^p_{HC}(X, \mathcal{H})$. Choisissons un hyperrecouvrement
$K$ de $X$ et un élément $\sigma \in \mathcal{H}(K_p)$ représentant
$\xi$ en cohomologie. Il existe une suite exacte correspondante de
complexes
$$
0 \to s(\mathcal{F}(K)) \to s(\mathcal{G}(K)) \to s(\mathcal{H}(K))
$$
mais rien ne garantit qu'il y ait aussi un zéro à droite, et c'est
la seule chose qui
nous empêche de définir $\delta(\xi)$ par une simple application du
lemme du serpent. Rappelons que
$$
\mathcal{H}(K_p) = \prod \mathcal{H}(U_i)
$$
si $K_p = \{U_i \to X\}$. Écrivons $\sigma =\prod \sigma_i$ avec
$\sigma_i \in \mathcal{H}(U_i)$. Comme $\mathcal{G} \to \mathcal{H}$ est
un morphisme surjectif de faisceaux, il existe des recouvrements
$\{U_{i, j} \to U_i\}$ tels que $\sigma_i|_{U_{i, j}}$ soit
l'image d'un élément $\tau_{i, j} \in \mathcal{G}(U_{i, j})$.
Considérons l'objet $Z = \{U_{i, j} \to X\}$ de la catégorie
$\text{SR}(\mathcal{C}, X)$ et son morphisme évident
$u : Z \to K_p$. Il est clair que $u$ est un recouvrement ; voir
la Définition \ref{definition-covering-SR}. D'après le
Lemme \ref{lemma-covering}, il
existe un morphisme $L \to K$ d'hyperrecouvrements de $X$ tel que
$L_p \to K_p$ se factorise par $u$. Après avoir remplacé $K$ par $L$,
on peut donc supposer que $\sigma$ est l'image d'un
élément $\tau \in \mathcal{G}(K_p)$. Remarquons que $d(\sigma) = 0$,
mais pas nécessairement que $d(\tau) = 0$. Ainsi, $d(\tau) \in \mathcal{F}(K_{p + 1})$
est un cocycle. Dans cette situation, on définit
$\delta(\xi)$ comme la classe du cocycle $d(\tau)$ dans
$\check{H}^{p + 1}_{HC}(X, \mathcal{F})$.

\medskip\noindent
Il faut maintenant vérifier plusieurs points :
(a) $\delta(\xi)$ ne dépend pas du choix de $\tau$,
(b) $\delta(\xi)$ ne dépend pas du choix de l'hyperrecouvrement
$L \to K$ sur lequel $\sigma$ se relève, et
(c) $\delta(\xi)$ ne dépend pas de l'hyperrecouvrement initial ni de
$\sigma$ choisi pour représenter $\xi$. Nous omettons la vérification de
(a), (b) et (c) ; l'indépendance des choix d'hyperrecouvrements
se ramène en fait aux Lemmes \ref{lemma-product-hypercoverings}
et \ref{lemma-homotopy}. Nous omettons aussi de vérifier que
$\delta$ est fonctoriel à l'égard des morphismes de suites exactes
courtes de faisceaux abéliens sur $\mathcal{C}$.

\medskip\noindent
Enfin, il faut vérifier qu'avec cette définition de $\delta$,
la suite exacte courte de faisceaux abéliens ci-dessus fournit une
suite exacte longue de groupes de cohomologie de {\v C}ech.
Montrons d'abord que si $\delta(\xi) = 0$ (avec $\xi$ comme ci-dessus), alors
$\xi$ est l'image d'un élément
$\xi' \in \check{H}^p_{HC}(X, \mathcal{G})$.
En effet, si $\delta(\xi) = 0$, alors, avec les notations précédentes, on
voit que la classe de $d(\tau)$ est nulle dans
$\check{H}^{p + 1}_{HC}(X, \mathcal{F})$. Il existe donc
un morphisme d'hyperrecouvrements $L \to K$ tel que la restriction
de $d(\tau)$ à un élément de $\mathcal{F}(L_{p + 1})$ soit
égale à $d(\upsilon)$ pour un certain $\upsilon \in \mathcal{F}(L_p)$.
Cela implique que $\tau|_{L_p} + \upsilon$ forme un
cocycle et détermine une classe $\xi' \in \check{H}^p(L, \mathcal{G})$
qui s'envoie sur $\xi$, comme voulu.

\medskip\noindent
Nous omettons la démonstration du fait que si $\xi' \in \check{H}^{p + 1}_{HC}(X, \mathcal{F})$
s'envoie sur zéro dans $\check{H}^{p + 1}_{HC}(X, \mathcal{G})$, alors il est
égal à $\delta(\xi)$ pour un certain $\xi \in \check{H}^p_{HC}(X, \mathcal{H})$.
\end{proof}

\noindent
Déduisons maintenant le cas de Verdier du
Théorème \ref{theorem-cohomology-hypercoverings}
par un artifice.

\begin{proposition}
\label{proposition-cohomology-hypercoverings}
Soit $\mathcal{C}$ un site admettant des produits fibrés et les produits de deux objets.
Soit $\mathcal{F}$ un faisceau abélien sur $\mathcal{C}$.
Soit $i \geq 0$. Alors
\begin{enumerate}
\item pour tout $\xi \in H^i(\mathcal{F})$, il existe un hyperrecouvrement
$K$ tel que $\xi$ appartienne à l'image de l'application canonique
$\check{H}^i(K, \mathcal{F}) \to H^i(\mathcal{F})$, et
\item si $K, L$ sont des hyperrecouvrements et $\xi_K \in \check{H}^i(K, \mathcal{F})$,
$\xi_L \in \check{H}^i(L, \mathcal{F})$ des éléments s'envoyant
sur le même élément de $H^i(\mathcal{F})$, alors il existe
un hyperrecouvrement $M$ et des morphismes $M \to K$ et $M \to L$ tels
que $\xi_K$ et $\xi_L$ s'envoient sur le même élément de
$\check{H}^i(M, \mathcal{F})$.
\end{enumerate}
Autrement dit, à des questions ensemblistes près, les groupes de
cohomologie de $\mathcal{F}$ sur $\mathcal{C}$ sont la limite inductive des
groupes de cohomologie de {\v C}ech de $\mathcal{F}$ sur tous les hyperrecouvrements.
\end{proposition}

\begin{proof}
Ce résultat est une conséquence immédiate du
Théorème \ref{theorem-cohomology-hypercoverings}.
En effet, on peut remplacer artificiellement $\mathcal{C}$ par un site
légèrement plus grand $\mathcal{C}'$ tel que
(I) $\mathcal{C}'$ possède un objet final $X$ et (II) les
hyperrecouvrements dans $\mathcal{C}$ soient essentiellement la
même chose que les hyperrecouvrements de $X$ dans $\mathcal{C}'$.
Mais la nature de la construction impose un assez grand nombre de
vérifications.

\medskip\noindent
Appelons une famille de morphismes $\{U_i \to U\}$ de $\mathcal{C}$
de but fixé un {\it recouvrement faible} si l'application
$\coprod_{i \in I} h_{U_i} \to h_U$ devient surjective après passage aux faisceaux associés.
Construisons un nouveau site $\mathcal{C}'$ comme suit :
\begin{enumerate}
\item comme catégorie, on pose $\Ob(\mathcal{C}') = \Ob(\mathcal{C}) \amalg \{X\}$
et on ajoute un unique morphisme vers $X$ depuis tout objet de $\mathcal{C}'$,
\item $\mathcal{C}'$ admet des produits fibrés dès que les produits fibrés et les produits
de deux objets existent dans $\mathcal{C}$,
\item les recouvrements de $\mathcal{C}'$ sont les recouvrements faibles de $\mathcal{C}$,
ainsi que les $\{U_i \to X\}_{i \in I}$ tels que soit $U_i = X$
pour un certain $i$, soit $U_i \not = X$ pour tout $i$ et que l'application
$\coprod h_{U_i} \to *$ de préfaisceaux sur $\mathcal{C}$ devienne
surjective après passage aux faisceaux associés sur $\mathcal{C}$,
\item on applique Ensembles, Lemme \ref{sets-lemma-coverings-site},
pour restreindre les recouvrements et obtenir notre site $\mathcal{C}'$.
\end{enumerate}
Alors $\Sh(\mathcal{C}') = \Sh(\mathcal{C})$, car le foncteur d'inclusion
$\mathcal{C} \to \mathcal{C}'$ est spécialement cocontinu
(voir Sites, Définition \ref{sites-definition-special-cocontinuous-functor}).
Nous omettons les vérifications immédiates.

\medskip\noindent
Choisissons un recouvrement $\{U_i \to X\}$ de $\mathcal{C}'$ tel que $U_i$ soit un
objet de $\mathcal{C}$ pour tout $i$ (cela est possible puisque
$\mathcal{C} \to \mathcal{C}'$ est spécialement cocontinu).
Alors $K_0 = \{U_i \to X\}$ est un recouvrement dans le
site $\mathcal{C}'$ construit ci-dessus. Considérons $K_0$ comme un objet de
$\text{SR}(\mathcal{C}', X)$ et posons $K_{init} = \text{cosk}_0(K_0)$.
Alors $K_{init}$ est un hyperrecouvrement de $X$ ; voir
l'Exemple \ref{example-cech}. Remarquons que chaque $K_{init, n}$ est de la forme
$\{W_j \to X\}$ avec $W_j \in \Ob(\mathcal{C})$.

\medskip\noindent
Démonstration de (1). Choisissons $\xi \in H^i(\mathcal{F}) = H^i(X, \mathcal{F}')$,
où $\mathcal{F}'$ est le faisceau abélien sur $\mathcal{C}'$ correspondant
à $\mathcal{F}$ sur $\mathcal{C}$. D'après le
Théorème \ref{theorem-cohomology-hypercoverings},
il existe un morphisme d'hyperrecouvrements $K' \to K_{init}$
de $X$ dans $\mathcal{C}'$ tel que $\xi$ provienne d'un élément
de $\check{H}^i(K', \mathcal{F})$.
Écrivons $K'_n = \{U_{n, j} \to X\}$. Comme $K'_n$ s'envoie dans
$K_{init, n}$, on voit que $U_{n, j}$ est un objet de $\mathcal{C}$.
On peut donc définir un objet simplicial $K$ de $\text{SR}(\mathcal{C})$
en posant $K_n = \{U_{n, j}\}$. Puisque les recouvrements de
$\mathcal{C}'$ formés de familles de morphismes de $\mathcal{C}$
sont des recouvrements faibles, on voit que $K$ est un hyperrecouvrement au sens
de la Définition \ref{definition-hypercovering-variant}.
Enfin, puisque $\mathcal{F}'$ est l'unique faisceau sur $\mathcal{C}'$
dont la restriction à $\mathcal{C}$ est égale à $\mathcal{F}$,
on voit que les complexes de {\v C}ech $s(\mathcal{F}(K))$
et $s(\mathcal{F}'(K'))$ sont identiques, ce qui démontre (1).
(La compatibilité avec l'application vers les groupes de cohomologie est omise.)

\medskip\noindent
Démonstration de (2). Soient $K$ et $L$ des hyperrecouvrements dans $\mathcal{C}$.
Soient $K'$ et $L'$ les objets simpliciaux de $\text{SR}(\mathcal{C}', X)$
obtenus à partir de $K$ et $L$ par le foncteur
$\text{SR}(\mathcal{C}) \to \text{SR}(\mathcal{C}', X)$,
$\{U_i\} \mapsto \{U_i \to X\}$. Comme précédemment, les complexes de
{\v C}ech sont égaux ; on obtient donc $\xi_{K'}$ et
$\xi_{L'}$ qui s'envoient sur la même classe de cohomologie de $\mathcal{F}'$
sur $\mathcal{C}'$. Après avoir éventuellement élargi notre choix
de recouvrements dans $\mathcal{C}'$ (en raison d'une question ensembliste),
on peut supposer que $K'$ et $L'$ sont des hyperrecouvrements de $X$ dans
$\mathcal{C}'$ ; cela résulte de notre définition des hyperrecouvrements dans la
Définition \ref{definition-hypercovering-variant} et
du fait que les recouvrements faibles de $\mathcal{C}$ donnent des recouvrements dans
$\mathcal{C}'$. D'après le
Théorème \ref{theorem-cohomology-hypercoverings},
il existe un hyperrecouvrement $M'$ de $X$ dans $\mathcal{C}'$
et des morphismes $M' \to K'$, $M' \to L'$ et $M' \to K_{init}$
tels que les restrictions de $\xi_{K'}$ et $\xi_{L'}$ soient le même élément de
$\check{H}^i(M', \mathcal{F})$. En explicitant cet énoncé comme ci-dessus,
on obtient (2).
\end{proof}





\section{Hyperrecouvrements d'espaces}
\label{section-hypercoverings-spaces}

\noindent
La théorie précédente présente déjà quelque intérêt dans le cas des espaces
topologiques. Dans ce cas, on peut expliciter ce qu'est un hyperrecouvrement et voir
ce que dit effectivement le résultat.

\medskip\noindent
Soit $X$ un espace topologique. Considérons le site $X_{Zar}$
de Sites, Exemple \ref{sites-example-site-topological}. Rappelons qu'un
objet de $X_{Zar}$ est simplement un ouvert de $X$ et que les morphismes
de $X_{Zar}$ correspondent simplement aux inclusions. Qu'est donc un
hyperrecouvrement de $X$ pour le site $X_{Zar}$ ?

\medskip\noindent
Explicitons d'abord la Définition \ref{definition-SR}.
Un objet de $\text{SR}(X_{Zar}, X)$ est simplement donné par un ensemble
$I$ et, pour chaque $i \in I$, un ouvert $U_i \subset X$.
Notons-le $\{U_i\}_{i \in I}$, puisqu'aucune
confusion n'est possible au sujet du morphisme $U_i \to X$.
Un morphisme $\{U_i\}_{i \in I} \to \{V_j\}_{j \in J}$
entre deux tels objets est donné par une application d'ensembles
$\alpha : I \to J$ telle que $U_i \subset V_{\alpha(i)}$ pour tout
$i \in I$. Quand un tel morphisme est-il un recouvrement ? C'est le cas
si et seulement si, pour tout $j \in J$, on a
$V_j = \bigcup_{i\in I, \ \alpha(i) = j} U_i$ (et que la famille est
un recouvrement dans le site $X_{Zar}$).

\medskip\noindent
Ce qui précède donne la description suivante d'un hyperrecouvrement
dans le site $X_{Zar}$. Un hyperrecouvrement de $X$ dans $X_{Zar}$
est donné par les données suivantes :
\begin{enumerate}
\item un ensemble simplicial $I$ (voir
Simplicial, section \ref{simplicial-section-simplicial-set}), et
\item pour tout $n \geq 0$ et tout $i \in I_n$, un ouvert $U_i \subset X$.
\end{enumerate}
Nous noterons une telle collection de données $(I, \{U_i\})$.
Pour que ce soit un hyperrecouvrement de $X$, on exige
les propriétés suivantes :
\begin{itemize}
\item pour $i \in I_n$ et $0 \leq a \leq n$,
on a $U_i \subset U_{d^n_a(i)}$,
\item pour $i \in I_n$ et $0 \leq a \leq n$, on a $U_i = U_{s^n_a(i)}$,
\item on a
\begin{equation}
\label{equation-covering-X}
X = \bigcup\nolimits_{i \in I_0} U_i,
\end{equation}
\item pour tous $i_0, i_1 \in I_0$, on a
\begin{equation}
\label{equation-covering-two}
U_{i_0} \cap U_{i_1} =
\bigcup\nolimits_{i \in I_1, \ d^1_0(i) = i_0, \ d^1_1(i) = i_1} U_i,
\end{equation}
\item pour tout $n \geq 1$ et tout
$(i_0, \ldots, i_{n + 1}) \in (I_n)^{n + 2}$ tels que
$d^n_{b - 1}(i_a) = d^n_a(i_b)$ pour tous $0\leq a < b\leq n + 1$,
on a
\begin{equation}
\label{equation-covering-general}
U_{i_0} \cap \ldots \cap U_{i_{n + 1}} =
\bigcup\nolimits_{i \in I_{n + 1},
\ d^{n + 1}_a(i) = i_a, \ a = 0, \ldots, n + 1} U_i,
\end{equation}
\item chacun des recouvrements ouverts (\ref{equation-covering-X}),
(\ref{equation-covering-two}) et (\ref{equation-covering-general})
est un élément de $\text{Cov}(X_{Zar})$
(c'est une condition ensembliste qui borne
la taille des ensembles d'indices des recouvrements).
\end{itemize}
Les conditions (\ref{equation-covering-X}) et
(\ref{equation-covering-two}) devraient être familières, par exemple d'après le
chapitre consacré aux faisceaux sur les espaces, et la condition
(\ref{equation-covering-general}) en est la généralisation naturelle.

\begin{remark}
\label{remark-not-covering-set}
Une particularité de cette description est que, si l'une des intersections
multiples $U_{i_0} \cap \ldots \cap U_{i_{n + 1}}$ est vide, alors
le recouvrement du membre de droite peut être le recouvrement vide.
Ainsi, les applications
$I_{n + 1} \to (\text{cosk}_n\text{sk}_n I)_{n + 1}$ ne sont pas nécessairement surjectives.
Cela signifie que la réalisation géométrique de $I$ peut être un espace
intéressant (non contractile).

\medskip\noindent
En effet, soit $I'_n \subset I_n$ le sous-ensemble
formé des simplexes $i \in I_n$ tels que
$U_i \not = \emptyset$. On vérifie aisément que $I' \subset I$
est un sous-ensemble simplicial et que $(I', \{U_i\})$ est un hyperrecouvrement.
On peut donc toujours raffiner un hyperrecouvrement en un hyperrecouvrement où
aucun des ouverts $U_i$ n'est vide.
\end{remark}

\begin{remark}
\label{remark-repackage-into-simplicial-space}
Reformulons encore une fois ces informations.
Supposons donc que $(I, \{U_i\})$ soit un hyperrecouvrement de
l'espace topologique $X$. À partir de ces données, on peut construire
un espace topologique simplicial $U_\bullet$ en posant
$$
U_n = \coprod\nolimits_{i \in I_n} U_i,
$$
et, pour $\varphi : [n] \to [m]$ donné, on définit
les morphismes $U(\varphi) : U_n \to U_m$ par les
inclusions $U_i \subset U_{\varphi(i)}$
pour $i \in I_n$. Cet espace topologique simplicial est
muni d'une augmentation $\epsilon : U_\bullet \to X$.
Avec ce morphisme, l'espace simplicial $U_\bullet$ devient
un hyperrecouvrement de $X$ le long duquel on a descente cohomologique
au sens de \cite[Expos\'e Vbis]{SGA4}.
Autrement dit, $H^n(U_\bullet, \epsilon^*\mathcal{F}) = H^n(X, \mathcal{F})$.
(Insérer ici une référence future à la cohomologie sur les espaces
simpliciaux et à la descente cohomologique formulée en ces termes.)
Supposons que $\mathcal{F}$ soit un faisceau abélien sur $X$.
Dans ce cas, la suite spectrale du Lemme \ref{lemma-cech-spectral-sequence}
devient la suite spectrale dont le terme $E_1$ est
$$
E_1^{p, q} = H^q(U_p, \epsilon_q^*\mathcal{F})
\Rightarrow
H^{p + q}(U_\bullet, \epsilon^*\mathcal{F}) = H^{p + q}(X, \mathcal{F})
$$
qui compare la cohomologie totale de $\epsilon^*\mathcal{F}$
aux groupes de cohomologie de $\mathcal{F}$ sur les morceaux
de $U_\bullet$. (Insérer une référence future à cette suite spectrale
ici.)
\end{remark}

\noindent
En topologie, on cherche souvent des hyperrecouvrements de $X$ qui
possèdent la propriété que tous les $U_i$ appartiennent à une base donnée de la topologie
de $X$ et que tous les recouvrements
(\ref{equation-covering-two}) et (\ref{equation-covering-general})
appartiennent à une collection cofinale donnée de recouvrements.
Voici deux lemmes qui en fournissent des exemples.

\begin{lemma}
\label{lemma-basis-hypercovering}
Soit $X$ un espace topologique.
Soit $\mathcal{B}$ une base de la topologie de $X$.
Il existe un hyperrecouvrement $(I, \{U_i\})$ de $X$
tel que chaque $U_i$ soit un élément de $\mathcal{B}$.
\end{lemma}

\begin{proof}
Soit $n \geq 0$.
Appelons {\it hyperrecouvrement $n$-tronqué} de $X$ une donnée
formée d'un ensemble simplicial $n$-tronqué $I$ et, pour chaque
$i \in I_a$, $0 \leq a \leq n$, d'un ouvert $U_i$ de $X$, de telle sorte que
les conditions définissant un hyperrecouvrement soient satisfaites lorsqu'elles ont un sens.
Autrement dit, on n'exige les relations d'inclusion et les conditions de
recouvrement que lorsque tous les simplexes qui y interviennent
sont des $a$-simplexes avec $a \leq n$. Le lemme découle de l'assertion
suivante : tout hyperrecouvrement $n$-tronqué $(I, \{U_i\})$ tel que
tous les $U_i \in \mathcal{B}$ se prolonge en un hyperrecouvrement $(n + 1)$-tronqué
sans ajout de $a$-simplexes pour $a \leq n$.
Procédons comme suit. Considérons d'abord l'ensemble simplicial
$(n + 1)$-tronqué $I'$ défini par
$I' = \text{sk}_{n + 1}(\text{cosk}_n I)$.
Rappelons que
$$
I'_{n + 1} =
\left\{
\begin{matrix}
(i_0, \ldots, i_{n + 1}) \in (I_n)^{n + 2} \text{ tels que}\\
d^n_{b - 1}(i_a) = d^n_a(i_b) \text{ pour tous }0\leq a < b\leq n + 1
\end{matrix}
\right\}
$$
Si $i' \in I'_{n + 1}$ est dégénéré, disons $i' = s^n_a(i)$, on pose
$U_{i'} = U_i$ (c'est de toute façon imposé par la deuxième condition).
On pose aussi $J_{i'} = \{i'\}$ dans ce cas.
Si $i' \in I'_{n + 1}$ est non dégénéré, disons
$i' = (i_0, \ldots, i_{n + 1})$, on choisit un ensemble
$J_{i'}$ et un recouvrement ouvert
\begin{equation}
\label{equation-choose-covering}
U_{i_0} \cap \ldots \cap U_{i_{n + 1}} =
\bigcup\nolimits_{i \in J_{i'}} U_i,
\end{equation}
avec $U_i \in \mathcal{B}$ pour $i \in J_{i'}$.
Posons
$$
I_{n + 1} = \coprod\nolimits_{i' \in I'_{n + 1}} J_{i'}
$$
Il existe une application canonique $\pi : I_{n + 1} \to I'_{n + 1}$ qui est
une bijection au-dessus de l'ensemble des simplexes dégénérés de $I'_{n + 1}$ par
construction.
Pour $i \in I_{n + 1}$, on définit $d^{n + 1}_a(i) = d^{n + 1}_a(\pi(i))$.
Pour $i \in I_n$, on définit $s^n_a(i) \in I_{n + 1}$ comme l'unique
simplexe situé au-dessus du simplexe dégénéré $s^n_a(i) \in I'_{n + 1}$.
Nous omettons la vérification que cela définit un hyperrecouvrement
$(n + 1)$-tronqué de $X$.
\end{proof}

\begin{lemma}
\label{lemma-quasi-separated-quasi-compact-hypercovering}
Soit $X$ un espace topologique.
Soit $\mathcal{B}$ une base de la topologie de $X$.
Supposons que
\begin{enumerate}
\item $X$ soit quasi-compact,
\item chaque $U \in \mathcal{B}$ soit un ouvert quasi-compact, et
\item l'intersection de deux ouverts quasi-compacts quelconques de
$X$ soit quasi-compacte.
\end{enumerate}
Alors il existe un hyperrecouvrement $(I, \{U_i\})$ de $X$ ayant les
propriétés suivantes :
\begin{enumerate}
\item chaque $U_i$ est un élément de la base $\mathcal{B}$,
\item chacun des $I_n$ est un ensemble fini, et, en particulier,
\item chacun des recouvrements (\ref{equation-covering-X}),
(\ref{equation-covering-two}) et (\ref{equation-covering-general})
est fini.
\end{enumerate}
\end{lemma}

\begin{proof}
Cela résulte directement de la construction donnée dans la démonstration du
Lemme \ref{lemma-basis-hypercovering}, si l'on choisit des recouvrements finis
par des éléments de $\mathcal{B}$ dans (\ref{equation-choose-covering}).
Nous omettons les détails.
\end{proof}







\section{Construction d'hyperrecouvrements}
\label{section-hypercovering-sites}

\noindent
Soit $\mathcal{C}$ un site. Dans cette section, nous considérerons un
objet simplicial de $\text{SR}(\mathcal{C})$ de la manière suivante.
Comme d'habitude, on pose $K_n = K([n])$ et on note $K(\varphi) : K_n \to K_m$
le morphisme associé à $\varphi : [m] \to [n]$.
On peut écrire $K_n = \{U_{n, i}\}_{i \in I_n}$. Pour
$\varphi : [m] \to [n]$, le morphisme $K(\varphi) : K_n \to K_m$
est donné par une application $\alpha(\varphi) : I_n \to I_m$ et des morphismes
$f_{\varphi, i} : U_{n, i} \to U_{m, \alpha(\varphi)(i)}$
pour $i \in I_n$. Le fait que $K$ soit un objet simplicial de
$\text{SR}(\mathcal{C})$ implique que $(I_n, \alpha(\varphi))$
est un ensemble simplicial
et que $f_{\psi, \alpha(\varphi)(i)} \circ f_{\varphi, i} =
f_{\varphi \circ \psi, i}$ lorsque $\psi : [l] \to [m]$.

\begin{lemma}
\label{lemma-split}
Soit $\mathcal{C}$ un site. Soit $K$ un objet simplicial $r$-tronqué
de $\text{SR}(\mathcal{C})$. Les conditions suivantes sont équivalentes :
\begin{enumerate}
\item $K$ est scindé (Simplicial, Définition \ref{simplicial-definition-split}),
\item $f_{\varphi, i} : U_{n, i} \to U_{m, \alpha(\varphi)(i)}$
est un isomorphisme pour $r \geq n \geq 0$,
$\varphi : [m] \to [n]$ surjective, $i \in I_n$, et
\item $f_{\sigma^n_j, i} : U_{n, i} \to U_{n + 1, \alpha(\sigma^n_j)(i)}$
est un isomorphisme pour $0 \leq j \leq n < r$, $i \in I_n$.
\end{enumerate}
Le même énoncé vaut pour les objets simpliciaux si, dans (2) et (3),
on pose $r = \infty$.
\end{lemma}

\begin{proof}
Le scindage d'un ensemble simplicial est unique et donné par
les indices non dégénérés $N(I_n)$ en chaque degré $n$ ; voir
Simplicial, Lemme \ref{simplicial-lemma-splitting-simplicial-sets}.
Le coproduit de deux objets $\{U_i\}_{i \in I}$ et $\{U_j\}_{j \in J}$
de $\text{SR}(\mathcal{C})$ est donné par $\{U_l\}_{l \in I \amalg J}$,
avec les notations évidentes. Un scindage de $K$ doit donc être donné par
$N(K_n) = \{U_i\}_{i \in N(I_n)}$. L'équivalence de (1) et (2)
résulte alors de l'explicitation des définitions. L'équivalence de (2)
et (3) résulte du fait que toute surjection
$\varphi : [m] \to [n]$ est un composé de morphismes
$\sigma^k_j$ avec $k = n, n + 1, \ldots, m - 1$.
\end{proof}

\begin{lemma}
\label{lemma-hypercovering-object}
Soit $\mathcal{C}$ un site admettant des produits fibrés.
Soit $\mathcal{B} \subset \Ob(\mathcal{C})$ un sous-ensemble.
Supposons que
\begin{enumerate}
\item tout objet $U$ de $\mathcal{C}$ admette un recouvrement
$\{U_j \to U\}_{j \in J}$ avec $U_j \in \mathcal{B}$, et
\item si $\{U_j \to U\}_{j \in J}$ est un recouvrement
avec $U_j \in \mathcal{B}$ et si $\{U' \to U\}$ est un morphisme avec
$U' \in \mathcal{B}$, alors $\{U_j \to U\}_{j \in J} \amalg \{U' \to U\}$
soit un recouvrement.
\end{enumerate}
Alors, pour tout $X$ de $\mathcal{C}$, il existe un hyperrecouvrement $K$
de $X$ tel que $K_n = \{U_{n, i}\}_{i \in I_n}$
avec $U_{n, i} \in \mathcal{B}$ pour tout $i \in I_n$.
\end{lemma}

\begin{proof}
La démonstration du Lemme \ref{lemma-basis-hypercovering}
constitue un échauffement pour celle-ci, et
nous encourageons le lecteur à la lire d'abord.

\medskip\noindent
On commence par remplacer $\mathcal{C}$ par le site $\mathcal{C}/X$.
Après ce remplacement, on peut supposer que $X$ est l'objet final
de $\mathcal{C}$ et que $\mathcal{C}$ admet toutes les limites finies
(Catégories, Lemme \ref{categories-lemma-finite-limits-exist}).

\medskip\noindent
Soit $n \geq 0$. Appelons
{\it hyperrecouvrement $n$-tronqué de type $\mathcal{B}$ de $X$}
un objet simplicial $n$-tronqué $K$
de $\text{SR}(\mathcal{C})$
tel que, pour $i \in I_a$, $0 \leq a \leq n$,
on ait $U_{a, i} \in \mathcal{B}$, que
$K_0$ soit un recouvrement de $X$ et que
$K_{a + 1} \to (\text{cosk}_a \text{sk}_a K)_{a + 1}$
soit un recouvrement pour $a = 0, \ldots, n - 1$,
comme dans la Définition \ref{definition-covering-SR}.

\medskip\noindent
Puisque, par hypothèse, $X$ admet un recouvrement $\{U_{0, i} \to X\}_{i \in I_0}$
avec $U_i \in \mathcal{B}$, on obtient un hyperrecouvrement
$0$-tronqué de type $\mathcal{B}$ de $X$. Observons que tout hyperrecouvrement
$0$-tronqué de type $\mathcal{B}$ de $X$ est scindé ; voir le
Lemme \ref{lemma-split}.

\medskip\noindent
Le lemme découle de l'assertion suivante, pour $n \geq 0$ : étant donné un
hyperrecouvrement $n$-tronqué de type $\mathcal{B}$, noté $K$, de $X$ et scindé,
on peut le prolonger en un
hyperrecouvrement $(n + 1)$-tronqué de type $\mathcal{B}$, scindé, de $X$.

\medskip\noindent
Construction du prolongement. Considérons l'objet simplicial $(n + 1)$-tronqué
$K' = \text{sk}_{n + 1}(\text{cosk}_n K)$ de $\text{SR}(\mathcal{C})$.
Écrivons
$$
K'_{n + 1} = \{U'_{n + 1, i}\}_{i \in I'_{n + 1}}
$$
Puisque $K = \text{sk}_n K'$, on a $K_a = K'_a$ pour $0 \leq a \leq n$.
Pour tout $i' \in I'_{n + 1}$, choisissons un recouvrement
\begin{equation}
\label{equation-choose-covering-B}
\{g_{n + 1, j} : U_{n + 1, j} \to U'_{n + 1, i'}\}_{j \in J_{i'}}
\end{equation}
avec $U_{n + 1, j} \in \mathcal{B}$ pour $j \in J_{i'}$.
C'est possible d'après l'hypothèse du lemme sur $\mathcal{B}$.
Pour $0 \leq m \leq n$, notons $N_m \subset I_m$ le sous-ensemble des
indices non dégénérés. Posons
$$
I_{n + 1} =
\coprod\nolimits_{\varphi : [n + 1] \to [m]\text{ surjective, }0\leq m \leq n}
N_m \amalg
\coprod\nolimits_{i' \in I'_{n + 1}} J_{i'}
$$
Pour $j \in I_{n + 1}$, posons
$$
U_{n + 1, j} =
\left\{
\begin{matrix}
U_{m, i} & \text{si} &
j = (\varphi, i) & \text{où} & \varphi : [n + 1] \to [m], i \in N_m \\
U_{n + 1, j} & \text{si} &
j \in J_{i'} & \text{où} & i' \in I'_{n + 1}
\end{matrix}
\right.
$$
avec les notations évidentes. Posons $K_{n + 1} = \{U_{n + 1, j}\}_{j \in I_{n + 1}}$.
Par construction, $U_{n + 1, j}$ est un élément
de $\mathcal{B}$ pour tout $j \in I_{n + 1}$. Définissons des
applications compatibles
$$
I_{n + 1} \to I'_{n + 1}
\quad\text{et}\quad
K_{n + 1} \to K'_{n + 1}
$$
La première application est donnée par
$(\varphi, i) \mapsto \alpha'(\varphi)(i)$ et
$(j \in J_{i'}) \mapsto i'$.
Pour la deuxième, on utilise les morphismes
$$
f'_{\varphi, i} : U_{m, i} \to U'_{n + 1, \alpha'(\varphi)(i)}
\quad\text{et}\quad
g_{n + 1, j} : U_{n + 1, j} \to U'_{n + 1, i'}
$$
Nous affirmons que le morphisme
$$
K_{n + 1} \to K'_{n + 1} =
(\text{cosk}_n \text{sk}_n K')_{n + 1} =
(\text{cosk}_n K)_{n + 1}
$$
est un recouvrement au sens de la Définition \ref{definition-covering-SR}.
En effet, si $i' \in I'_{n + 1}$, alors soit $i'$ est non dégénéré
et l'image inverse de $i'$ dans $I_{n + 1}$ est égale à $J_{i'}$,
auquel cas on obtient un recouvrement de $U'_{n + 1, i'}$ par notre choix
(\ref{equation-choose-covering-B}), soit $i'$ est dégénéré et
l'image inverse de $i'$ dans $I_{n + 1}$ est
$J_{i'} \amalg \{(\varphi, i)\}$ pour un unique couple $(\varphi, i)$,
auquel cas notre choix (\ref{equation-choose-covering-B})
et l'hypothèse (2) du lemme donnent un recouvrement.

\medskip\noindent
Pour achever la démonstration, il faut définir les morphismes
$K(\varphi) : K_{n + 1} \to K_m$ correspondant aux morphismes
$\varphi : [m] \to [n + 1]$, $0 \leq m \leq n$, ainsi que les morphismes
$K(\varphi) : K_m \to K_{n + 1}$ correspondant aux morphismes
$\varphi : [n + 1] \to [m]$, $0 \leq m \leq n$,
de manière à satisfaire les relations de composition convenables.
Pour le premier type, on utilise le composé
$$
K_{n + 1} \to K'_{n + 1} \xrightarrow{K'(\varphi)} K'_m = K_m
$$
pour définir $K(\varphi) : K_{n + 1} \to K_m$.
Pour le second type, supposons donné $\varphi : [n + 1] \to [m]$,
$0 \leq m \leq n$. On définit le morphisme correspondant
$K(\varphi) : K_m \to K_{n + 1}$ comme suit :
\begin{enumerate}
\item pour $i \in I_m$, il existe une unique application surjective
$\psi : [m] \to [m_0]$ et un unique $i_0 \in I_{m_0}$ non dégénéré
tels que $\alpha(\psi)(i_0) = i$\footnote{Par exemple, si $i$ est
non dégénéré, alors $m = m_0$ et $\psi = \text{id}_{[m]}$.},
\item on pose $\varphi_0 = \psi_0 \circ \varphi : [n + 1] \to [m_0]$
et on envoie
$i \in I_m$ sur $(\varphi_0, i_0) \in I_{n + 1}$ ; autrement dit,
$\alpha(\varphi)(i) = (\varphi_0, i_0)$, et
\item le morphisme
$f_{\varphi, i} : U_{m, i} \to U_{n + 1, \alpha(\varphi)(i)} = U_{m_0, i_0}$
est l'inverse de l'isomorphisme $f_{\psi, i_0} : U_{m_0, i_0} \to U_{m, i}$
(voir le Lemme \ref{lemma-split}).
\end{enumerate}
Nous omettons la vérification immédiate mais laborieuse que cela définit
un hyperrecouvrement $(n + 1)$-tronqué de type $\mathcal{B}$, scindé, de $X$
prolongeant celui qui était donné en degré $n$. Tout est en fait clair
d'après ce qui précède, sauf la vérification que les morphismes
$K(\varphi)$ se composent correctement pour tout $\varphi : [a] \to [b]$
avec $0 \leq a, b \leq n + 1$.
\end{proof}

\begin{lemma}
\label{lemma-hypercovering-site}
Soit $\mathcal{C}$ un site admettant des égalisateurs et des produits fibrés.
Soit $\mathcal{B} \subset \Ob(\mathcal{C})$ un sous-ensemble. Supposons
que tout objet de $\mathcal{C}$ possède un recouvrement
dont les membres sont des éléments de $\mathcal{B}$.
Alors il existe un hyperrecouvrement $K$ tel que
$K_n = \{U_i\}_{i \in I_n}$ avec $U_i \in \mathcal{B}$
pour tout $i \in I_n$.
\end{lemma}

\begin{proof}
Cette démonstration est presque identique à celle du
Lemme \ref{lemma-hypercovering-object}. Nous n'en
expliquerons que les différences.

\medskip\noindent
Soit $n \geq 1$. Appelons
{\it hyperrecouvrement $n$-tronqué de type $\mathcal{B}$}
un objet simplicial
$n$-tronqué $K$ de $\text{SR}(\mathcal{C})$
tel que, pour $i \in I_a$, $0 \leq a \leq n$,
on ait $U_{a, i} \in \mathcal{B}$ et que
\begin{enumerate}
\item $F(K_0)^\# \to *$ soit surjectif,
\item $F(K_1)^\# \to F(K_0)^\# \times F(K_0)^\#$ soit surjectif,
\item $F(K_{a + 1})^\# \to F((\text{cosk}_a \text{sk}_a K)_{a + 1})^\#$
soit surjectif pour $a = 1, \ldots, n - 1$.
\end{enumerate}
Construisons d'abord explicitement un hyperrecouvrement $1$-tronqué de type $\mathcal{B}$ et scindé.

\medskip\noindent
Prenons $I_0 = \mathcal{B}$ et $K_0 = \{U\}_{U \in \mathcal{B}}$.
Alors (1) est vérifiée d'après notre hypothèse sur $\mathcal{B}$. Posons
$$
\Omega =
\{(U, V, W, a, b) \mid U, V, W \in \mathcal{B}, a : U \to V, b : U \to W\}
$$
Puis posons $I_1 = I_0 \amalg \Omega$. Pour $i \in I_1$, on pose
$U_{1, i} = U_{0, i}$ si $i \in I_0$, et $U_{1, i} = U$
si $i = (U, V, W, a, b) \in \Omega$. L'application
$K(\sigma^0_0) : K_0 \to K_1$ correspond à
l'inclusion $\alpha(\sigma^0_0) : I_0 \to I_1$
et à l'identité $f_{\sigma^0_0, i} : U_{0, i} \to U_{1, i}$
sur les objets. Les applications $K(\delta^1_0), K(\delta^1_1) : K_1 \to K_0$
correspondent aux deux applications $I_1 \to I_0$ qui sont
l'identité sur $I_0 \subset I_1$ et envoient $(U, V, W, a, b) \in \Omega \subset I_1$
sur $V$, resp.\ $W$. Les morphismes correspondants
$f_{\delta^1_0, i}, f_{\delta^1_1, i} : U_{1, i} \to U_{0, i}$ sont
l'identité si $i \in I_0$, et $a, b$ si $i = (U, V, W, a, b) \in \Omega$.
La condition (2) est vérifiée parce que toute section de
$F(K_0)^\# \times F(K_0)^\#$ sur un objet $U$ de $\mathcal{C}$
provient, après remplacement de $U$ par les membres d'un recouvrement,
d'une application $U \to F(K_0) \times F(K_0)$.
Cela signifie à son tour qu'on dispose de $V, W \in \mathcal{B}$
et de deux morphismes $U \to V$ et $U \to W$. En remplaçant encore
$U$ par les membres d'un recouvrement, on peut supposer que $U \in \mathcal{B}$,
comme voulu.

\medskip\noindent
Le lemme découle de l'assertion suivante : étant donné un
hyperrecouvrement $n$-tronqué de type $\mathcal{B}$, noté $K$ et scindé, pour $n \geq 1$,
on peut le prolonger en un hyperrecouvrement $(n + 1)$-tronqué de type $\mathcal{B}$ et scindé.
Ici, l'argument procède exactement comme dans la démonstration du
Lemme \ref{lemma-hypercovering-object}.
Nous omettons les détails précis, à l'exception des remarques suivantes.
Premièrement, l'hypothèse (2) de la démonstration du présent
lemme est inutile, car nous n'avons pas besoin que le morphisme
$K_{n + 1} \to (\text{cosk}_n K)_{n + 1}$ soit un recouvrement ;
il suffit qu'il induise une surjection sur les faisceaux d'ensembles associés,
ce qui résulte de
Sites, Lemme \ref{sites-lemma-covering-surjective-after-sheafification}.
Deuxièmement, l'hypothèse selon laquelle $\mathcal{C}$ admet des produits fibrés et des égalisateurs
garantit que $\text{SR}(\mathcal{C})$ admet des produits fibrés
et des égalisateurs, et que $F$ commute à ceux-ci
(Lemme \ref{lemma-coprod-prod-SR}). Cela suffit à
garantir l'existence des foncteurs cosquelette employés (voir
Simplicial, Remarque \ref{simplicial-remark-existence-cosk} et
Catégories, Lemme \ref{categories-lemma-fibre-products-equalizers-exist}).
\end{proof}

\begin{lemma}
\label{lemma-hypercovering-morphism-sites}
Soit $f : \mathcal{C} \to \mathcal{D}$ un morphisme de sites
donné par le foncteur $u : \mathcal{D} \to \mathcal{C}$.
Supposons que $\mathcal{D}$ et $\mathcal{C}$ admettent des égalisateurs et des
produits fibrés, et que $u$ commute à ceux-ci.
Si un objet simplicial $K$ de $\text{SR}(\mathcal{D})$
est un hyperrecouvrement, alors $u(K)$ est un hyperrecouvrement.
\end{lemma}

\begin{proof}
Si l'on écrit $K_n = \{U_{n, i}\}_{i \in I_n}$ comme dans l'introduction
de cette section, alors $u(K)$ est l'objet de $\text{SR}(\mathcal{C})$
donné par $u(K_n) = \{u(U_i)\}_{i \in I_n}$.
D'après Sites, Lemme \ref{sites-lemma-pullback-representable-sheaf},
on a $f^{-1}h_U^\# = h_{u(U)}^\#$ pour $U \in \Ob(\mathcal{D})$.
Cela signifie que $f^{-1}F(K_n)^\# = F(u(K_n))^\#$ pour tout $n$.
Vérifions les conditions (1), (2) et (3) pour que $u(K)$ soit un
hyperrecouvrement selon la Définition \ref{definition-hypercovering-variant}.
Comme $f^{-1}$ est un foncteur exact, on trouve que
$$
F(u(K_0))^\# = f^{-1}F(K_0)^\# \to f^{-1}* = *
$$
est surjectif comme image inverse d'un morphisme surjectif, ce qui donne (1).
De même,
$$
F(u(K_1))^\# = f^{-1}F(K_1)^\# \to
f^{-1} (F(K_0) \times F(K_0))^\# = F(u(K_0))^\# \times F(u(K_0))^\#
$$
est surjectif comme image inverse, ce qui donne (2). Pour la condition (3),
il suffit, afin de conclure par la même méthode, que
$$
F((\text{cosk}_n \text{sk}_n u(K))_{n + 1})^\# =
f^{-1}F((\text{cosk}_n \text{sk}_n K)_{n + 1})^\#
$$
Ce qui précède montre que $f^{-1}F(-) = F(u(-))$. Il suffit donc de montrer
que $u$ commute aux limites utilisées pour définir
$(\text{cosk}_n \text{sk}_n K)_{n + 1}$ pour $n \geq 1$.
D'après Simplicial, Remarque \ref{simplicial-remark-existence-cosk},
ces limites sont des limites finies connexes, auxquelles $u$ commute
par hypothèse.
\end{proof}

\begin{lemma}
\label{lemma-hypercovering-continuous-functor}
Soient $\mathcal{C}$, $\mathcal{D}$ des sites. Soit
$u : \mathcal{D} \to \mathcal{C}$ un foncteur continu.
Supposons que $\mathcal{D}$ et $\mathcal{C}$ admettent des produits fibrés
et que $u$ commute à ceux-ci. Soient $Y \in \mathcal{D}$ et
$K \in \text{SR}(\mathcal{D}, Y)$ un hyperrecouvrement de $Y$.
Alors $u(K)$ est un hyperrecouvrement de $u(Y)$.
\end{lemma}

\begin{proof}
C'est plus facile que la démonstration du Lemme \ref{lemma-hypercovering-morphism-sites},
car la notion d'hyperrecouvrement d'un objet est plus forte ; voir les
Définitions \ref{definition-hypercovering} et \ref{definition-covering-SR}.
En effet, $u$ envoie les recouvrements sur des recouvrements par définition
d'un morphisme de sites. Il suffit de vérifier que $u$ commute aux
limites utilisées pour définir
$(\text{cosk}_n \text{sk}_n K)_{n + 1}$ pour $n \geq 1$.
C'est clair, car le foncteur induit
$\mathcal{D}/Y \to \mathcal{C}/X$ commute à toutes les limites finies
(et la source et le but admettent toutes les limites finies d'après
Catégories, Lemme \ref{categories-lemma-finite-limits-exist}).
\end{proof}

\begin{lemma}
\label{lemma-w-contractible}
Soit $\mathcal{C}$ un site. Soit $\mathcal{B} \subset \Ob(\mathcal{C})$
un sous-ensemble. Supposons que
\begin{enumerate}
\item $\mathcal{C}$ admette des produits fibrés,
\item pour tout $X \in \Ob(\mathcal{C})$, il existe un recouvrement fini
$\{U_i \to X\}_{i \in I}$ avec $U_i \in \mathcal{B}$,
\item si $\{U_i \to X\}_{i \in I}$ est un recouvrement fini avec
$U_i \in \mathcal{B}$ et si $U \to X$ est un morphisme avec $U \in \mathcal{B}$,
alors $\{U_i \to X\}_{i \in I} \amalg \{U \to X\}$ soit un recouvrement.
\end{enumerate}
Alors, pour tout $X$, il existe un hyperrecouvrement $K$ de $X$
tel que chaque $K_n = \{U_{n, i} \to X\}_{i \in I_n}$, où
$I_n$ est fini et $U_{n, i} \in \mathcal{B}$.
\end{lemma}

\begin{proof}
Ce lemme est l'analogue du
Lemme \ref{lemma-quasi-separated-quasi-compact-hypercovering}
pour les sites. Pour le démontrer, on suit exactement la démonstration du
Lemme \ref{lemma-hypercovering-object},
en prêtant attention aux deux points suivants :
\begin{enumerate}
\item[(a)] on choisit le recouvrement initial $\{U_{0, i} \to X\}_{i \in I_0}$
avec $U_{0, i} \in \mathcal{B}$ de telle sorte que l'ensemble d'indices $I_0$ soit fini, et
\item[(b)] lors du choix des recouvrements
(\ref{equation-choose-covering-B}),
on choisit $J_{i'}$ fini.
\end{enumerate}
Le lecteur voit aisément qu'avec ces modifications, on obtient
des ensembles d'indices finis $I_n$ pour tout $n$.
\end{proof}

\begin{remark}
\label{remark-taking-disjoint-unions}
Soit $\mathcal{C}$ un site. Soient
$K$ et $L$ des objets de $\text{SR}(\mathcal{C})$.
Écrivons $K = \{U_i\}_{i \in I}$ et $L = \{V_j\}_{j \in J}$.
Supposons que $U = \coprod_{i \in I} U_i$ et $V = \coprod_{j \in J} V_j$
existent. On obtient alors
$$
\Mor_{\text{SR}(\mathcal{C})}(K, L) \longrightarrow \Mor_\mathcal{C}(U, V)
$$
comme suit. Étant donné $f : K \to L$, défini par $\alpha : I \to J$
et les $f_i : U_i \to V_{\alpha(i)}$, on obtient une transformation de foncteurs
$$
\Mor_\mathcal{C}(V, -) =
\prod\nolimits_{j \in J} \Mor_\mathcal{C}(V_j, -)
\to
\prod\nolimits_{i \in I} \Mor_\mathcal{C}(U_i, -) =
\Mor_\mathcal{C}(U, -)
$$
qui envoie $(g_j)_{j \in J}$ sur
$(g_{\alpha(i)} \circ f_i)_{i \in I}$. Le lemme de Yoneda
produit donc l'application correspondante $U \to V$. Bien entendu, $U \to V$
envoie le facteur $U_i$ dans le facteur $V_{\alpha(i)}$ par
le morphisme $f_i$.
\end{remark}

\begin{remark}
\label{remark-take-unions-hypercovering}
Soit $\mathcal{C}$ un site. Supposons que $\mathcal{C}$ admette
des produits fibrés et des égalisateurs, et soit $K$ un hyperrecouvrement.
Écrivons $K_n = \{U_{n, i}\}_{i \in I_n}$. Supposons que
\begin{enumerate}
\item[(a)] $U_n = \coprod_{i \in I_n} U_{n, i}$ existe, et
\item[(b)] $\coprod_{i \in I_n} h_{U_{n, i}} \to h_{U_n}$ induise
un isomorphisme sur les faisceaux associés.
\end{enumerate}
On obtient alors un autre objet simplicial $L$ de $\text{SR}(\mathcal{C})$
avec $L_n = \{U_n\}$ ; voir
la Remarque \ref{remark-taking-disjoint-unions}.
Nous affirmons que $L$ est un hyperrecouvrement.
Vérifions pour cela les conditions (1), (2) et (3) de la
Définition \ref{definition-hypercovering-variant}.
La condition (1) résulte de (b) et de la condition (1) pour $K$.
La condition (2) se démontre exactement de la même manière.
La condition (3) résulte des égalités
\begin{align*}
F((\text{cosk}_n \text{sk}_n L)_{n + 1})^\#
& =
((\text{cosk}_n \text{sk}_n F(L)^\#)_{n + 1}) \\
& =
((\text{cosk}_n \text{sk}_n F(K)^\#)_{n + 1}) \\
& =
F((\text{cosk}_n \text{sk}_n K)_{n + 1})^\#
\end{align*}
pour $n \geq 1$ ; la condition pour $K$ implique donc celle pour
$L$, exactement comme dans (1) et (2).
Remarquons que $F$ commute aux limites connexes et que le foncteur faisceau associé est exact,
ce qui démontre la première et la dernière égalité ; l'égalité du milieu résulte de
$F(K)^\# = F(L)^\#$ d'après (b).
\end{remark}

\begin{remark}
\label{remark-take-unions-hypercovering-X}
Soit $\mathcal{C}$ un site. Soit $X \in \Ob(\mathcal{C})$.
Supposons que $\mathcal{C}$ admette des produits fibrés, et soit $K$ un hyperrecouvrement de $X$.
Écrivons $K_n = \{U_{n, i}\}_{i \in I_n}$. Supposons que
\begin{enumerate}
\item[(a)] $U_n = \coprod_{i \in I_n} U_{n, i}$ existe,
\item[(b)] étant donnés des morphismes
$(\alpha, f_i) : \{U_i\}_{i \in I} \to \{V_j\}_{j \in J}$ et
$(\beta, g_k) : \{W_k\}_{k \in K} \to \{V_j\}_{j \in J}$
dans $\text{SR}(\mathcal{C})$, tels que
$U = \coprod U_i$, $V = \coprod V_j$ et $W = \coprod W_j$
existent, alors $U \times_V W =
\coprod_{(i, j, k), \alpha(i) = j = \beta(k)} U_i \times_{V_j} W_k$,
\item[(c)] si $(\alpha, f_i) : \{U_i\}_{i \in I} \to \{V_j\}_{j \in J}$
est un recouvrement au sens de la
Définition \ref{definition-covering-SR}
et si $U = \coprod U_i$ et $V = \coprod V_j$ existent,
alors le morphisme correspondant $U \to V$
de la Remarque \ref{remark-taking-disjoint-unions}
est un recouvrement de $\mathcal{C}$.
\end{enumerate}
On obtient alors un autre objet simplicial $L$ de $\text{SR}(\mathcal{C})$
avec $L_n = \{U_n\}$ ; voir
la Remarque \ref{remark-taking-disjoint-unions}.
Nous affirmons que $L$ est un hyperrecouvrement de $X$.
Vérifions pour cela les conditions (1) et (2) de la
Définition \ref{definition-hypercovering}.
La condition (1) résulte de (c) et de la condition (1) pour $K$,
car la condition (1) pour $K$ dit que $K_0 = \{U_{0, i}\}_{i \in I_0}$
est un recouvrement de $\{X\}$ au sens de la
Définition \ref{definition-covering-SR}.
La condition (2) résulte du fait que $\mathcal{C}/X$ admet
toutes les limites finies, donc que $\text{SR}(\mathcal{C}/X)$
admet toutes les limites finies, et du fait que la condition (b) dit que la
construction « prendre les sommes disjointes » commute
à ces limites finies. Ainsi, le morphisme
$$
L_{n + 1} \longrightarrow (\text{cosk}_n \text{sk}_n L)_{n + 1}
$$
est un recouvrement, puisqu'il résulte de l'application de notre
foncteur « prendre les sommes disjointes » au morphisme
$$
K_{n + 1} \longrightarrow (\text{cosk}_n \text{sk}_n K)_{n + 1}
$$
qui est supposé être un recouvrement au sens de la
Définition \ref{definition-covering-SR}, d'après la condition (2) pour $K$.
Cela a un sens parce que la propriété (b) garantit en particulier
que, si l'on part d'un diagramme fini d'objets
semi-représentables au-dessus de $X$
pour lesquels on peut prendre les sommes disjointes, alors
la limite du diagramme dans $\text{SR}(\mathcal{C}/X)$
est encore un objet semi-représentable au-dessus de $X$ pour lequel
on peut prendre les sommes disjointes.
\end{remark}





\begin{multicols}{2}[\section{Autres chapitres}]
\noindent
Préliminaires
\begin{enumerate}
\item \hyperref[introduction-section-phantom]{Introduction}
\item \hyperref[conventions-section-phantom]{Conventions}
\item \hyperref[sets-section-phantom]{Théorie des ensembles}
\item \hyperref[categories-section-phantom]{Catégories}
\item \hyperref[topology-section-phantom]{Topologie}
\item \hyperref[sheaves-section-phantom]{Faisceaux sur les espaces}
\item \hyperref[sites-section-phantom]{Sites et faisceaux}
\item \hyperref[stacks-section-phantom]{Champs}
\item \hyperref[fields-section-phantom]{Corps}
\item \hyperref[algebra-section-phantom]{Algèbre commutative}
\item \hyperref[brauer-section-phantom]{Groupes de Brauer}
\item \hyperref[homology-section-phantom]{Algèbre homologique}
\item \hyperref[derived-section-phantom]{Catégories dérivées}
\item \hyperref[simplicial-section-phantom]{Méthodes simpliciales}
\item \hyperref[more-algebra-section-phantom]{Compléments d'algèbre}
\item \hyperref[smoothing-section-phantom]{Lissage des morphismes d'anneaux}
\item \hyperref[modules-section-phantom]{Faisceaux de modules}
\item \hyperref[sites-modules-section-phantom]{Modules sur les sites}
\item \hyperref[injectives-section-phantom]{Objets injectifs}
\item \hyperref[cohomology-section-phantom]{Cohomologie des faisceaux}
\item \hyperref[sites-cohomology-section-phantom]{Cohomologie sur les sites}
\item \hyperref[dga-section-phantom]{Algèbre différentielle graduée}
\item \hyperref[dpa-section-phantom]{Algèbre à puissances divisées}
\item \hyperref[sdga-section-phantom]{Faisceaux différentiels gradués}
\item \hyperref[hypercovering-section-phantom]{Hyperrecouvrements}
\end{enumerate}
Schémas
\begin{enumerate}
\setcounter{enumi}{25}
\item \hyperref[schemes-section-phantom]{Schémas}
\item \hyperref[constructions-section-phantom]{Constructions de schémas}
\item \hyperref[properties-section-phantom]{Propriétés des schémas}
\item \hyperref[morphisms-section-phantom]{Morphismes de schémas}
\item \hyperref[coherent-section-phantom]{Cohomologie des schémas}
\item \hyperref[divisors-section-phantom]{Diviseurs}
\item \hyperref[limits-section-phantom]{Limites de schémas}
\item \hyperref[varieties-section-phantom]{Variétés}
\item \hyperref[topologies-section-phantom]{Topologies sur les schémas}
\item \hyperref[descent-section-phantom]{Descente}
\item \hyperref[perfect-section-phantom]{Catégories dérivées de schémas}
\item \hyperref[more-morphisms-section-phantom]{Compléments sur les morphismes}
\item \hyperref[flat-section-phantom]{Compléments sur la platitude}
\item \hyperref[groupoids-section-phantom]{Schémas en groupoïdes}
\item \hyperref[more-groupoids-section-phantom]{Compléments sur les schémas en groupoïdes}
\item \hyperref[etale-section-phantom]{Morphismes étales de schémas}
\end{enumerate}
Thèmes de théorie des schémas
\begin{enumerate}
\setcounter{enumi}{41}
\item \hyperref[chow-section-phantom]{Homologie de Chow}
\item \hyperref[intersection-section-phantom]{Théorie de l'intersection}
\item \hyperref[pic-section-phantom]{Schémas de Picard des courbes}
\item \hyperref[weil-section-phantom]{Théories cohomologiques de Weil}
\item \hyperref[adequate-section-phantom]{Modules adéquats}
\item \hyperref[dualizing-section-phantom]{Complexes dualisants}
\item \hyperref[duality-section-phantom]{Dualité pour les schémas}
\item \hyperref[discriminant-section-phantom]{Discriminants et différentes}
\item \hyperref[derham-section-phantom]{Cohomologie de de Rham}
\item \hyperref[local-cohomology-section-phantom]{Cohomologie locale}
\item \hyperref[algebraization-section-phantom]{Géométrie algébrique et formelle}
\item \hyperref[curves-section-phantom]{Courbes algébriques}
\item \hyperref[resolve-section-phantom]{Résolution des surfaces}
\item \hyperref[models-section-phantom]{Réduction semi-stable}
\item \hyperref[functors-section-phantom]{Foncteurs et morphismes}
\item \hyperref[equiv-section-phantom]{Catégories dérivées de variétés}
\item \hyperref[pione-section-phantom]{Groupes fondamentaux de schémas}
\item \hyperref[etale-cohomology-section-phantom]{Cohomologie étale}
\item \hyperref[crystalline-section-phantom]{Cohomologie cristalline}
\item \hyperref[proetale-section-phantom]{Cohomologie pro-étale}
\item \hyperref[relative-cycles-section-phantom]{Cycles relatifs}
\item \hyperref[more-etale-section-phantom]{Compléments de cohomologie étale}
\item \hyperref[trace-section-phantom]{La formule des traces}
\end{enumerate}
Espaces algébriques
\begin{enumerate}
\setcounter{enumi}{64}
\item \hyperref[spaces-section-phantom]{Espaces algébriques}
\item \hyperref[spaces-properties-section-phantom]{Propriétés des espaces algébriques}
\item \hyperref[spaces-morphisms-section-phantom]{Morphismes d'espaces algébriques}
\item \hyperref[decent-spaces-section-phantom]{Espaces algébriques décents}
\item \hyperref[spaces-cohomology-section-phantom]{Cohomologie des espaces algébriques}
\item \hyperref[spaces-limits-section-phantom]{Limites d'espaces algébriques}
\item \hyperref[spaces-divisors-section-phantom]{Diviseurs sur les espaces algébriques}
\item \hyperref[spaces-over-fields-section-phantom]{Espaces algébriques sur des corps}
\item \hyperref[spaces-topologies-section-phantom]{Topologies sur les espaces algébriques}
\item \hyperref[spaces-descent-section-phantom]{Descente et espaces algébriques}
\item \hyperref[spaces-perfect-section-phantom]{Catégories dérivées d'espaces}
\item \hyperref[spaces-more-morphisms-section-phantom]{Compléments sur les morphismes d'espaces}
\item \hyperref[spaces-flat-section-phantom]{Platitude sur les espaces algébriques}
\item \hyperref[spaces-groupoids-section-phantom]{Groupoïdes dans les espaces algébriques}
\item \hyperref[spaces-more-groupoids-section-phantom]{Compléments sur les groupoïdes dans les espaces}
\item \hyperref[bootstrap-section-phantom]{Amorçage}
\item \hyperref[spaces-pushouts-section-phantom]{Sommes amalgamées d'espaces algébriques}
\end{enumerate}
Thèmes de géométrie
\begin{enumerate}
\setcounter{enumi}{81}
\item \hyperref[spaces-chow-section-phantom]{Groupes de Chow des espaces}
\item \hyperref[groupoids-quotients-section-phantom]{Quotients de groupoïdes}
\item \hyperref[spaces-more-cohomology-section-phantom]{Compléments sur la cohomologie des espaces}
\item \hyperref[spaces-simplicial-section-phantom]{Espaces simpliciaux}
\item \hyperref[spaces-duality-section-phantom]{Dualité pour les espaces}
\item \hyperref[formal-spaces-section-phantom]{Espaces algébriques formels}
\item \hyperref[restricted-section-phantom]{Algébrisation des espaces formels}
\item \hyperref[spaces-resolve-section-phantom]{Résolution des surfaces revisitée}
\end{enumerate}
Théorie des déformations
\begin{enumerate}
\setcounter{enumi}{89}
\item \hyperref[formal-defos-section-phantom]{Théorie formelle des déformations}
\item \hyperref[defos-section-phantom]{Théorie des déformations}
\item \hyperref[cotangent-section-phantom]{Le complexe cotangent}
\item \hyperref[examples-defos-section-phantom]{Problèmes de déformation}
\end{enumerate}
Champs algébriques
\begin{enumerate}
\setcounter{enumi}{93}
\item \hyperref[algebraic-section-phantom]{Champs algébriques}
\item \hyperref[examples-stacks-section-phantom]{Exemples de champs}
\item \hyperref[stacks-sheaves-section-phantom]{Faisceaux sur les champs algébriques}
\item \hyperref[criteria-section-phantom]{Critères de représentabilité}
\item \hyperref[artin-section-phantom]{Axiomes d'Artin}
\item \hyperref[quot-section-phantom]{Espaces de Quot et de Hilbert}
\item \hyperref[stacks-properties-section-phantom]{Propriétés des champs algébriques}
\item \hyperref[stacks-morphisms-section-phantom]{Morphismes de champs algébriques}
\item \hyperref[stacks-limits-section-phantom]{Limites de champs algébriques}
\item \hyperref[stacks-cohomology-section-phantom]{Cohomologie des champs algébriques}
\item \hyperref[stacks-perfect-section-phantom]{Catégories dérivées de champs}
\item \hyperref[stacks-introduction-section-phantom]{Introduction aux champs algébriques}
\item \hyperref[stacks-more-morphisms-section-phantom]{Compléments sur les morphismes de champs}
\item \hyperref[stacks-geometry-section-phantom]{La géométrie des champs}
\end{enumerate}
Thèmes de théorie des espaces de modules
\begin{enumerate}
\setcounter{enumi}{107}
\item \hyperref[moduli-section-phantom]{Champs de modules}
\item \hyperref[moduli-curves-section-phantom]{Espaces de modules de courbes}
\end{enumerate}
Divers
\begin{enumerate}
\setcounter{enumi}{109}
\item \hyperref[examples-section-phantom]{Exemples}
\item \hyperref[exercises-section-phantom]{Exercices}
\item \hyperref[guide-section-phantom]{Guide bibliographique}
\item \hyperref[desirables-section-phantom]{Desiderata}
\item \hyperref[coding-section-phantom]{Style de programmation}
\item \hyperref[obsolete-section-phantom]{Obsolète}
\item \hyperref[fdl-section-phantom]{Licence de documentation libre GNU}
\item \hyperref[index-section-phantom]{Index généré automatiquement}
\end{enumerate}
\end{multicols}


\bibliography{my}
\bibliographystyle{amsalpha}

\end{document}
